% Môn: Vật lý
% Lớp: 12
% Chương: Chương II. Khí lí tưởng
% Bài: L12C2 Bài 9. Định luật Boyle
% Số câu: 160
% App đọc file này trực tiếp — sửa rồi Commit trên GitHub, bấm Đồng bộ GitHub .tex

% L12C2 Bài 9. Định luật Boyle

% ===== Câu 1 =====
% ID: L12C2B9-01-DS
% Mức: TH
\dangbt{Bài toán tổng hợp về quá trình đẳng nhiệt}
\begin{ex}
% ID: L12C2B9-01-DS
% Mức: TH
Xét các phát biểu sau trong dạng “Bài toán tổng hợp về quá trình đẳng nhiệt”.
\choiceTF
{\True Đại lượng nào không đổi trong quá trình đẳng nhiệt của một lượng khí xác định? — phát biểu đúng là: Nhiệt độ.}
{Đại lượng nào không đổi trong quá trình đẳng nhiệt của một lượng khí xác định? — phát biểu Áp suất là đúng.}
{\True Hệ thức nào là định luật Boyle? — phát biểu đúng là: $p_1V_1=p_2V_2$.}
{Hệ thức nào là định luật Boyle? — phát biểu $V_1/T_1=V_2/T_2$ là đúng.}
\loigiai{
\begin{itemchoice}
\item (Đúng) Đại lượng nào không đổi trong quá trình đẳng nhiệt của một lượng khí xác định? — phát biểu đúng là: Nhiệt độ. (Vì): Trong quá trình đẳng nhiệt, nhiệt độ của lượng khí xác định là không đổi
\item (Sai) Đại lượng nào không đổi trong quá trình đẳng nhiệt của một lượng khí xác định? — phát biểu Áp suất là đúng. (Vì): Trong quá trình đẳng nhiệt, áp suất của lượng khí xác định không phải lúc nào cũng không đổi; nó thay đổi khi thể tích thay đổi theo định luật Boyle
\item (Đúng) Hệ thức nào là định luật Boyle? — phát biểu đúng là: $p_1V_1=p_2V_2$. (Vì): ịnh luật Boyle phát biểu rằng tích của áp suất và thể tích của một lượng khí xác định là không đổi ở nhiệt độ không đổi, tức là $p_1V_1=p_2V_2$
\item (Sai) Hệ thức nào là định luật Boyle? — phát biểu $V_1/T_1=V_2/T_2$ là đúng. (Vì): Hệ thức $V_1/T_1=V_2/T_2$ mô tả định luật Charles (hoặc Gay-Lussac thứ nhất), áp dụng cho quá trình đẳng áp, không phải định luật Boyle
\end{itemchoice}
}
\end{ex}

% ===== Câu 2 =====
% ID: L12C2B9-01-TL
% Mức: VD
\begin{bt}
% ID: L12C2B9-01-TL
% Mức: VD
Trong dạng “Bài toán tổng hợp về quá trình đẳng nhiệt”, Thể tích giảm từ 28 lít còn 7 lít ở nhiệt độ không đổi. Áp suất tăng bao nhiêu lần?
\loigiai{
$p_2/p_1=V_1/V_2=4$.
}
\end{bt}

% ===== Câu 3 =====
% ID: L12C2B9-01-TLN
% Mức: TH
\begin{ex}
% ID: L12C2B9-01-TLN
% Mức: TH
Trong dạng “Bài toán tổng hợp về quá trình đẳng nhiệt”, Khí 6 lít ở $100\,\text{kPa}$ được nén đẳng nhiệt còn 3 lít. Tính áp suất sau theo kPa.
\shortans{200}
\loigiai{
$p_2=100\cdot6/3=200$ kPa.
}
\end{ex}

% ===== Câu 4 =====
% ID: L12C2B9-01-TN
% Mức: NB
\begin{ex}
% ID: L12C2B9-01-TN
% Mức: NB
Trong dạng “Bài toán tổng hợp về quá trình đẳng nhiệt”, Đại lượng nào không đổi trong quá trình đẳng nhiệt của một lượng khí xác định?
\choice
{\True Nhiệt độ}
{Áp suất}
{Thể tích}
{Khối lượng riêng}
\loigiai{
Đáp án A. Quá trình đẳng nhiệt có nhiệt độ không đổi.
}
\end{ex}

% ===== Câu 5 =====
% ID: L12C2B9-02-DS
% Mức: TH
\begin{ex}
% ID: L12C2B9-02-DS
% Mức: TH
Xét các phát biểu sau trong dạng “Bài toán tổng hợp về quá trình đẳng nhiệt”.
\choiceTF
{\True Khí có áp suất $130\,\text{kPa}$, thể tích 10 lít được nén đẳng nhiệt còn 5 lít. Áp suất sau là — phát biểu đúng là: $260\,\text{kPa}$.}
{Khí có áp suất $130\,\text{kPa}$, thể tích 10 lít được nén đẳng nhiệt còn 5 lít. Áp suất sau là — phát biểu $65\,\text{kPa}$ là đúng.}
{\True Khí có áp suất $90\,\text{kPa}$, thể tích 7 lít giãn đẳng nhiệt đến 21 lít. Áp suất sau là — phát biểu đúng là: $30\,\text{kPa}$.}
{Khí có áp suất $90\,\text{kPa}$, thể tích 7 lít giãn đẳng nhiệt đến 21 lít. Áp suất sau là — phát biểu $90\,\text{kPa}$ là đúng.}
\loigiai{
\begin{itemchoice}
\item (Đúng) Khí có áp suất $130\,\text{kPa}$, thể tích 10 lít được nén đẳng nhiệt còn 5 lít. Áp suất sau là — phát biểu đúng là: $260\,\text{kPa}$. (Vì): Trong quá trình đẳng nhiệt, tích áp suất và thể tích không đổi: $p_1V_1 = p_2V_2$. Với $p_1 = 130\,\text{kPa}$, $V_1 = 10\,\text{lít}$, $V_2 = 5\,\text{lít}$, ta có $p_2 = \frac{p_1V_1}{V_2} = \frac{130 \times 10}{5} = 260\,\text{kPa}$
\item (Sai) Khí có áp suất $130\,\text{kPa}$, thể tích 10 lít được nén đẳng nhiệt còn 5 lít. Áp suất sau là — phát biểu $65\,\text{kPa}$ là đúng. (Vì): Phát biểu này cho rằng áp suất sau là $65\,\text{kPa}$, điều này mâu thuẫn với kết quả tính toán ở mệnh đề A
\item (Đúng) Khí có áp suất $90\,\text{kPa}$, thể tích 7 lít giãn đẳng nhiệt đến 21 lít. Áp suất sau là — phát biểu đúng là: $30\,\text{kPa}$. (Vì): Trong quá trình đẳng nhiệt, tích áp suất và thể tích không đổi: $p_1V_1 = p_2V_2$. Với $p_1 = 90\,\text{kPa}$, $V_1 = 7\,\text{lít}$, $V_2 = 21\,\text{lít}$, ta có $p_2 = \frac{p_1V_1}{V_2} = \frac{90 \times 7}{21} = 30\,\text{kPa}$
\item (Sai) Khí có áp suất $90\,\text{kPa}$, thể tích 7 lít giãn đẳng nhiệt đến 21 lít. Áp suất sau là — phát biểu $90\,\text{kPa}$ là đúng. (Vì): Phát biểu này cho rằng áp suất sau là $90\,\text{kPa}$, điều này mâu thuẫn với kết quả tính toán ở mệnh đề C
\end{itemchoice}
}
\end{ex}

% ===== Câu 6 =====
% ID: L12C2B9-02-TL
% Mức: VD
\begin{bt}
% ID: L12C2B9-02-TL
% Mức: VD
Trong dạng “Bài toán tổng hợp về quá trình đẳng nhiệt”, 130
\loigiai{
$p=390/3=130$ kPa.
}
\end{bt}

% ===== Câu 7 =====
% ID: L12C2B9-02-TLN
% Mức: TH
\begin{ex}
% ID: L12C2B9-02-TLN
% Mức: TH
Trong dạng “Bài toán tổng hợp về quá trình đẳng nhiệt”, Khí 5 lít ở $330\,\text{kPa}$ giãn đẳng nhiệt đến áp suất $110\,\text{kPa}$. Tính thể tích sau theo lít.
\shortans{15}
\loigiai{
$V_2=330\cdot5/110=15$ lít.
}
\end{ex}

% ===== Câu 8 =====
% ID: L12C2B9-02-TN
% Mức: NB
\begin{ex}
% ID: L12C2B9-02-TN
% Mức: NB
Trong dạng “Bài toán tổng hợp về quá trình đẳng nhiệt”, Hệ thức nào là định luật Boyle?
\choice
{\True $p_1V_1=p_2V_2$}
{$p_1/T_1=p_2/T_2$}
{$V_1/T_1=V_2/T_2$}
{$p_1V_2=p_2V_1$}
\loigiai{
Đáp án A. Ở nhiệt độ không đổi, $pV$ là hằng số.
}
\end{ex}

% ===== Câu 9 =====
% ID: L12C2B9-03-DS
% Mức: TH
\begin{ex}
% ID: L12C2B9-03-DS
% Mức: TH
Trong dạng “Bài toán tổng hợp về quá trình đẳng nhiệt”, Trình bày cơ sở lí thuyết của dạng “Bài toán tổng hợp về quá trình đẳng nhiệt” và nêu điều kiện áp dụng.
\choiceTF
{\True Trong quá trình đẳng nhiệt, tích của áp suất và thể tích của một lượng khí nhất định là không đổi.}
{Định luật Boyle áp dụng cho mọi loại khí.}
{Khi áp dụng định luật Boyle, có thể sử dụng áp suất tuyệt đối hoặc áp suất tương đối đều được.}
{Nhiệt độ tuyệt đối trong quá trình đẳng nhiệt có thể thay đổi.}
\loigiai{
\begin{itemchoice}
\item (Đúng) Trong quá trình đẳng nhiệt, tích của áp suất và thể tích của một lượng khí nhất định là không đổi. (Vì): ây là nội dung của định luật Boyle: $p_1V_1 = p_2V_2 = const$ khi nhiệt độ $T$ không đổi
\item (Sai) Định luật Boyle áp dụng cho mọi loại khí. (Vì): Định luật Boyle chỉ áp dụng cho khí lí tưởng
\item (Sai) Khi áp dụng định luật Boyle, có thể sử dụng áp suất tuyệt đối hoặc áp suất tương đối đều được. (Vì): Khi áp dụng định luật Boyle, phải sử dụng áp suất tuyệt đối
\item (Sai) Nhiệt độ tuyệt đối trong quá trình đẳng nhiệt có thể thay đổi. (Vì): Quá trình đẳng nhiệt là quá trình có nhiệt độ không đổi
\end{itemchoice}
}
\end{ex}

% ===== Câu 10 =====
% ID: L12C2B9-03-TL
% Mức: VD
\begin{bt}
% ID: L12C2B9-03-TL
% Mức: VD
Trong dạng “Bài toán tổng hợp về quá trình đẳng nhiệt”, $\rho_2/\rho_1=p_2/p_1=3$.
\loigiai{
$\rho_2/\rho_1=p_2/p_1=3$.
}
\end{bt}

% ===== Câu 11 =====
% ID: L12C2B9-03-TLN
% Mức: VD
\begin{ex}
% ID: L12C2B9-03-TLN
% Mức: VD
Trong dạng “Bài toán tổng hợp về quá trình đẳng nhiệt”, Thể tích giảm từ 28 lít còn 7 lít ở nhiệt độ không đổi. Áp suất tăng bao nhiêu lần?
\shortans{4}
\loigiai{
$p_2/p_1=V_1/V_2=4$.
}
\end{ex}

% ===== Câu 12 =====
% ID: L12C2B9-03-TN
% Mức: TH
\begin{ex}
% ID: L12C2B9-03-TN
% Mức: TH
Trong dạng “Bài toán tổng hợp về quá trình đẳng nhiệt”, Khí có áp suất $130\,\text{kPa}$, thể tích 10 lít được nén đẳng nhiệt còn 5 lít. Áp suất sau là
\choice
{\True $260\,\text{kPa}$}
{$65\,\text{kPa}$}
{$130\,\text{kPa}$}
{$520\,\text{kPa}$}
\loigiai{
Đáp án A. $p_2=p_1V_1/V_2=130\cdot10/5=260$ kPa.
}
\end{ex}

% ===== Câu 13 =====
% ID: L12C2B9-04-DS
% Mức: TH
\begin{ex}
% ID: L12C2B9-04-DS
% Mức: TH
Trong dạng “Bài toán tổng hợp về quá trình đẳng nhiệt”, Mô tả dạng đồ thị đặc trưng và cách nhận biết quá trình trong dạng “Bài toán tổng hợp về quá trình đẳng nhiệt”.
\choiceTF
{\True Đồ thị biểu diễn quá trình đẳng nhiệt trong hệ tọa độ (p, V) là một đường hypebol.}
{Đồ thị biểu diễn quá trình đẳng nhiệt trong hệ tọa độ (p, T) là một đường thẳng song song với trục p.}
{Đồ thị biểu diễn quá trình đẳng nhiệt trong hệ tọa độ ($V$, T) là một đường thẳng đi qua gốc tọa độ.}
{Để nhận biết quá trình đẳng nhiệt, ta cần kiểm tra xem tích pV có không đổi hay không.}
\loigiai{
\begin{itemchoice}
\item (Đúng) Đồ thị biểu diễn quá trình đẳng nhiệt trong hệ tọa độ (p, V) là một đường hypebol. (Vì): Theo định luật Boyle-Mariotte, với $T$ không đổi, ta có $pV = const$. Đồ thị $p(V)$ là một đường hypebol
\item (Sai) Đồ thị biểu diễn quá trình đẳng nhiệt trong hệ tọa độ (p, T) là một đường thẳng song song với trục p. (Vì): Đồ thị biểu diễn quá trình đẳng nhiệt trong hệ tọa độ (p, T) là một đường thẳng song song với trục $T$, vì p tỉ lệ nghịch với 1/ $T$, không phải song song với trục p
\item (Sai) Đồ thị biểu diễn quá trình đẳng nhiệt trong hệ tọa độ ($V$, T) là một đường thẳng đi qua gốc tọa độ. (Vì): onst)T$. Đồ thị $V(T)$ là một đường thẳng đi qua gốc tọa độ
\item (Sai) Để nhận biết quá trình đẳng nhiệt, ta cần kiểm tra xem tích pV có không đổi hay không.
\end{itemchoice}
}
\end{ex}

% ===== Câu 14 =====
% ID: L12C2B9-04-TL
% Mức: VDC
\begin{bt}
% ID: L12C2B9-04-TL
% Mức: VDC
Trong dạng “Bài toán tổng hợp về quá trình đẳng nhiệt”, undefined
\loigiai{
$V_2=150\cdot28/(600)=7$ lít.
}
\end{bt}

% ===== Câu 15 =====
% ID: L12C2B9-04-TLN
% Mức: VD
\begin{ex}
% ID: L12C2B9-04-TLN
% Mức: VD
Trong dạng “Bài toán tổng hợp về quá trình đẳng nhiệt”, Một lượng khí có $pV=390$ kPa.lít. Khi V=3 lít, tính p theo kPa.
\shortans{130}
\loigiai{
$p=390/3=130$ kPa.
}
\end{ex}

% ===== Câu 16 =====
% ID: L12C2B9-04-TN
% Mức: TH
\begin{ex}
% ID: L12C2B9-04-TN
% Mức: TH
Trong dạng “Bài toán tổng hợp về quá trình đẳng nhiệt”, Khí có áp suất $90\,\text{kPa}$, thể tích 7 lít giãn đẳng nhiệt đến 21 lít. Áp suất sau là
\choice
{\True $30\,\text{kPa}$}
{$270\,\text{kPa}$}
{$90\,\text{kPa}$}
{$10\,\text{kPa}$}
\loigiai{
Đáp án A. $p_2=90\cdot7/21=30$ kPa.
}
\end{ex}

% ===== Câu 17 =====
% ID: L12C2B9-05-DS
% Mức: VD
\begin{ex}
% ID: L12C2B9-05-DS
% Mức: VD
Xét các phát biểu sau trong dạng “Bài toán tổng hợp về quá trình đẳng nhiệt”.
\choiceTF
{\True Thể tích khí giảm 4 lần trong quá trình đẳng nhiệt thì áp suất — phát biểu đúng là: tăng 4 lần.}
{Thể tích khí giảm 4 lần trong quá trình đẳng nhiệt thì áp suất — phát biểu giảm 4 lần là đúng.}
{\True Trên hệ trục p- $V$, đường đẳng nhiệt là — phát biểu đúng là: một nhánh hypebol.}
{Trên hệ trục p- $V$, đường đẳng nhiệt là — phát biểu đường thẳng song song trục p là đúng.}
\loigiai{
\begin{itemchoice}
\item (Đúng) Thể tích khí giảm 4 lần trong quá trình đẳng nhiệt thì áp suất — phát biểu đúng là: tăng 4 lần. (Vì): Trong quá trình đẳng nhiệt, áp suất tỉ lệ nghịch với thể tích ($p \sim \frac{1}{V}$). Nếu thể tích giảm 4 lần thì áp suất tăng 4 lần
\item (Sai) Thể tích khí giảm 4 lần trong quá trình đẳng nhiệt thì áp suất — phát biểu giảm 4 lần là đúng. (Vì): Phát biểu này mâu thuẫn với định luật Boyle-Mariotte cho quá trình đẳng nhiệt
\item (Đúng) Trên hệ trục p- $V$, đường đẳng nhiệt là — phát biểu đúng là: một nhánh hypebol. (Vì): Phương trình trạng thái khí lí tưởng cho quá trình đẳng nhiệt là $pV = \text{hằng số}$. Trên hệ trục tọa độ $p-V$, đồ thị này là một nhánh của đường hypebol
\item (Sai) Trên hệ trục p- $V$, đường đẳng nhiệt là — phát biểu đường thẳng song song trục p là đúng. (Vì): Đường đẳng nhiệt trên hệ trục $p-V$ là một nhánh hypebol, không phải là đường thẳng song song với trục $p$
\end{itemchoice}
}
\end{ex}

% ===== Câu 18 =====
% ID: L12C2B9-05-TL
% Mức: TH
\dangbt{Khối khí chịu áp suất của cột chất lỏng}
\begin{ex}
% ID: L12C2B9-05-TL
% Mức: TH
Trong dạng “Khối khí chịu áp suất của cột chất lỏng”, Trình bày cơ sở lí thuyết của dạng “Khối khí chịu áp suất của cột chất lỏng” và nêu điều kiện áp dụng.
\choiceTF

\loigiai{
Nêu đúng đại lượng không đổi, hệ thức đặc trưng và yêu cầu dùng nhiệt độ kelvin, áp suất tuyệt đối khi có liên quan.
}
\end{ex}

% ===== Câu 19 =====
% ID: L12C2B9-05-TLN
% Mức: VD
\dangbt{Bài toán tổng hợp về quá trình đẳng nhiệt}
\begin{ex}
% ID: L12C2B9-05-TLN
% Mức: VD
Trong dạng “Bài toán tổng hợp về quá trình đẳng nhiệt”, Áp suất đẳng nhiệt tăng từ $140\,\text{kPa}$ lên $420\,\text{kPa}$. Khối lượng riêng tăng bao nhiêu lần?
\shortans{3}
\loigiai{
$\rho_2/\rho_1=p_2/p_1=3$.
}
\end{ex}

% ===== Câu 20 =====
% ID: L12C2B9-05-TN
% Mức: TH
\begin{ex}
% ID: L12C2B9-05-TN
% Mức: TH
Trong dạng “Bài toán tổng hợp về quá trình đẳng nhiệt”, Thể tích khí giảm 4 lần trong quá trình đẳng nhiệt thì áp suất
\choice
{\True tăng 4 lần}
{giảm 4 lần}
{không đổi}
{tăng 16 lần}
\loigiai{
Đáp án A. Do $p$ tỉ lệ nghịch với $V$.
}
\end{ex}

% ===== Câu 21 =====
% ID: L12C2B9-06-DS
% Mức: VD
\begin{ex}
% ID: L12C2B9-06-DS
% Mức: VD
Xét các phát biểu sau trong dạng “Bài toán tổng hợp về quá trình đẳng nhiệt”.
\choiceTF
{\True Khi dùng định luật Boyle phải sử dụng — phát biểu đúng là: áp suất tuyệt đối.}
{Khi dùng định luật Boyle phải sử dụng — phát biểu áp suất dư là đúng.}
{\True Muốn nén khí gần đẳng nhiệt nên thực hiện — phát biểu đúng là: chậm để khí trao đổi nhiệt.}
{Muốn nén khí gần đẳng nhiệt nên thực hiện — phát biểu ở thể tích không đổi là đúng.}
\loigiai{
\begin{itemchoice}
\item (Đúng) Khi dùng định luật Boyle phải sử dụng — phát biểu đúng là: áp suất tuyệt đối. (Vì): ịnh luật Boyle (hay định luật Mariotte) phát biểu rằng với một lượng khí xác định, ở nhiệt độ không đổi, tích của áp suất tuyệt đối và thể tích là một hằng số ($p V = \text{const}$). Do đó, khi áp dụng định luật này, ta phải sử dụng áp suất tuyệt đối
\item (Sai) Khi dùng định luật Boyle phải sử dụng — phát biểu áp suất dư là đúng. (Vì): Áp suất dư (hay áp suất tương đối) là hiệu số giữa áp suất tuyệt đối và áp suất khí quyển. Định luật Boyle áp dụng cho áp suất tuyệt đối, không phải áp suất dư
\item (Đúng) Muốn nén khí gần đẳng nhiệt nên thực hiện — phát biểu đúng là: chậm để khí trao đổi nhiệt. (Vì): Quá trình đẳng nhiệt là quá trình có nhiệt độ không đổi. Để nén khí gần như đẳng nhiệt, quá trình nén phải diễn ra chậm để khí có đủ thời gian trao đổi nhiệt với môi trường, giữ cho nhiệt độ không tăng lên đáng kể
\item (Sai) Muốn nén khí gần đẳng nhiệt nên thực hiện — phát biểu ở thể tích không đổi là đúng. (Vì): Quá trình đẳng nhiệt liên quan đến sự thay đổi của áp suất và thể tích khi nhiệt độ không đổi. Nếu thể tích không đổi, đó là quá trình đẳng tích, không phải đẳng nhiệt
\end{itemchoice}
}
\end{ex}

% ===== Câu 22 =====
% ID: L12C2B9-06-TL
% Mức: TH
\dangbt{Khối khí chịu áp suất của cột chất lỏng}
\begin{ex}
% ID: L12C2B9-06-TL
% Mức: TH
Trong dạng “Khối khí chịu áp suất của cột chất lỏng”, Mô tả dạng đồ thị đặc trưng và cách nhận biết quá trình trong dạng “Khối khí chịu áp suất của cột chất lỏng”.
\choiceTF

\loigiai{
Xác định các trục, đại lượng không đổi và suy ra hình dạng đồ thị từ hệ thức vật lí tương ứng.
}
\end{ex}

% ===== Câu 23 =====
% ID: L12C2B9-06-TLN
% Mức: VDC
\dangbt{Bài toán tổng hợp về quá trình đẳng nhiệt}
\begin{ex}
% ID: L12C2B9-06-TLN
% Mức: VDC
Trong dạng “Bài toán tổng hợp về quá trình đẳng nhiệt”, Khí 28 lít ở $150\,\text{kPa}$ giãn/nén đẳng nhiệt đến $600\,\text{kPa}$. Tính thể tích theo lít.
\shortans{7}
\loigiai{
$V_2=150\cdot28/(600)=7$ lít.
}
\end{ex}

% ===== Câu 24 =====
% ID: L12C2B9-06-TN
% Mức: VD
\begin{ex}
% ID: L12C2B9-06-TN
% Mức: VD
Trong dạng “Bài toán tổng hợp về quá trình đẳng nhiệt”, Trên hệ trục p- $V$, đường đẳng nhiệt là
\choice
{\True một nhánh hypebol}
{đường thẳng qua gốc}
{đường thẳng song song trục p}
{parabol}
\loigiai{
Đáp án A. Từ $p=C/V$, đồ thị là một nhánh hypebol.
}
\end{ex}

% ===== Câu 25 =====
% ID: L12C2B9-07-DS
% Mức: VDC
\begin{ex}
% ID: L12C2B9-07-DS
% Mức: VDC
Xét các phát biểu sau trong dạng “Bài toán tổng hợp về quá trình đẳng nhiệt”.
\choiceTF
{\True Với lượng khí cố định đẳng nhiệt, khối lượng riêng tỉ lệ — phát biểu đúng là: thuận với áp suất.}
{Với lượng khí cố định đẳng nhiệt, khối lượng riêng tỉ lệ — phát biểu nghịch với áp suất là đúng.}
{\True Trong ống có cột chất lỏng, áp suất khí bị giam được xác định bằng — phát biểu đúng là: cân bằng áp suất tại cùng độ cao.}
{Trong ống có cột chất lỏng, áp suất khí bị giam được xác định bằng — phát biểu chỉ $\rho gh$ là đúng.}
\loigiai{
\begin{itemchoice}
\item (Đúng) Với lượng khí cố định đẳng nhiệt, khối lượng riêng tỉ lệ — phát biểu đúng là: thuận với áp suất. (Vì): Với lượng khí cố định và quá trình đẳng nhiệt, ta có $pV = \text{hằng số}$. Khối lượng riêng $\rho = \frac{m}{V}$. Từ $V = \frac{m}{\rho}$, thay vào biểu thức trên ta được $p\frac{m}{\rho} = \text{hằng số}$. Do $m$ không đổi, suy ra $\frac{p}{\rho} = \text{hằng số}$, hay $\rho$ tỉ lệ thuận với $p$
\item (Sai) Với lượng khí cố định đẳng nhiệt, khối lượng riêng tỉ lệ — phát biểu nghịch với áp suất là đúng. (Vì): Phát biểu này trái với hệ thức đã được chứng minh ở mệnh đề A
\item (Đúng) Trong ống có cột chất lỏng, áp suất khí bị giam được xác định bằng — phát biểu đúng là: cân bằng áp suất tại cùng độ cao. (Vì): Áp suất khí bị giam trong ống có cột chất lỏng được xác định dựa trên nguyên lý cân bằng áp suất tại cùng một độ cao trong lòng chất lỏng
\item (Sai) Trong ống có cột chất lỏng, áp suất khí bị giam được xác định bằng — phát biểu chỉ $\rho gh$ là đúng. (Vì): Áp suất khí bị giam trong ống có cột chất lỏng được xác định bằng biểu thức $p_{\text{khí}} = p_0 \pm \rho gh$, trong đó $p_0$ là áp suất khí quyển, $\rho$ là khối lượng riêng của chất lỏng, $g$ là gia tốc trọng trường và $h$ là độ chênh lệch mực chất lỏng. Biểu thức chỉ $\rho gh$ không đủ để xác định áp suất khí bị giam
\end{itemchoice}
}
\end{ex}

% ===== Câu 26 =====
% ID: L12C2B9-07-TL
% Mức: VD
\dangbt{Khối khí chịu áp suất của cột chất lỏng}
\begin{bt}
% ID: L12C2B9-07-TL
% Mức: VD
Trong dạng “Khối khí chịu áp suất của cột chất lỏng”, Thể tích giảm từ 10 lít còn 5 lít ở nhiệt độ không đổi. Áp suất tăng bao nhiêu lần?
\loigiai{
$p_2/p_1=V_1/V_2=2$.
}
\end{bt}

% ===== Câu 27 =====
% ID: L12C2B9-07-TLN
% Mức: TH
\begin{ex}
% ID: L12C2B9-07-TLN
% Mức: TH
Trong dạng “Khối khí chịu áp suất của cột chất lỏng”, Khí 21 lít ở $140\,\text{kPa}$ được nén đẳng nhiệt còn 7 lít. Tính áp suất sau theo kPa.
\shortans{420}
\loigiai{
$p_2=140\cdot21/7=420$ kPa.
}
\end{ex}

% ===== Câu 28 =====
% ID: L12C2B9-07-TN
% Mức: VD
\dangbt{Bài toán tổng hợp về quá trình đẳng nhiệt}
\begin{ex}
% ID: L12C2B9-07-TN
% Mức: VD
Trong dạng “Bài toán tổng hợp về quá trình đẳng nhiệt”, Khi dùng định luật Boyle phải sử dụng
\choice
{\True áp suất tuyệt đối}
{áp suất dư}
{áp suất khí quyển}
{chênh lệch áp suất bất kì}
\loigiai{
Đáp án A. Phương trình trạng thái dùng áp suất tuyệt đối của khí.
}
\end{ex}

% ===== Câu 29 =====
% ID: L12C2B9-08-DS
% Mức: TH
\dangbt{Khối khí chịu áp suất của cột chất lỏng}
\begin{ex}
% ID: L12C2B9-08-DS
% Mức: TH
Xét các phát biểu sau trong dạng “Khối khí chịu áp suất của cột chất lỏng”.
\choiceTF
{\True Đại lượng nào không đổi trong quá trình đẳng nhiệt của một lượng khí xác định? — phát biểu đúng là: Nhiệt độ.}
{Đại lượng nào không đổi trong quá trình đẳng nhiệt của một lượng khí xác định? — phát biểu Áp suất là đúng.}
{\True Hệ thức nào là định luật Boyle? — phát biểu đúng là: $p_1V_1=p_2V_2$.}
{Hệ thức nào là định luật Boyle? — phát biểu $V_1/T_1=V_2/T_2$ là đúng.}
\loigiai{
\begin{itemchoice}
\item (Đúng) Đại lượng nào không đổi trong quá trình đẳng nhiệt của một lượng khí xác định? — phát biểu đúng là: Nhiệt độ. (Vì): Quá trình đẳng nhiệt có nhiệt độ không đổi. b) Sai: trái với hệ thức hoặc điều kiện đã nêu. c) Đúng: Ở nhiệt độ không đổi, $pV$ là hằng số. d) Sai: trái với hệ thức hoặc điều kiện đã nêu
\item (Sai) Đại lượng nào không đổi trong quá trình đẳng nhiệt của một lượng khí xác định? — phát biểu Áp suất là đúng.
\item (Đúng) Hệ thức nào là định luật Boyle? — phát biểu đúng là: $p_1V_1=p_2V_2$.
\item (Sai) Hệ thức nào là định luật Boyle? — phát biểu $V_1/T_1=V_2/T_2$ là đúng.
\end{itemchoice}
}
\end{ex}

% ===== Câu 30 =====
% ID: L12C2B9-08-TL
% Mức: VD
\begin{bt}
% ID: L12C2B9-08-TL
% Mức: VD
Trong dạng “Khối khí chịu áp suất của cột chất lỏng”, 110
\loigiai{
$p=770/7=110$ kPa.
}
\end{bt}

% ===== Câu 31 =====
% ID: L12C2B9-08-TLN
% Mức: TH
\begin{ex}
% ID: L12C2B9-08-TLN
% Mức: TH
Trong dạng “Khối khí chịu áp suất của cột chất lỏng”, Khí 3 lít ở $600\,\text{kPa}$ giãn đẳng nhiệt đến áp suất $150\,\text{kPa}$. Tính thể tích sau theo lít.
\shortans{12}
\loigiai{
$V_2=600\cdot3/150=12$ lít.
}
\end{ex}

% ===== Câu 32 =====
% ID: L12C2B9-08-TN
% Mức: VD
\dangbt{Bài toán tổng hợp về quá trình đẳng nhiệt}
\begin{ex}
% ID: L12C2B9-08-TN
% Mức: VD
Trong dạng “Bài toán tổng hợp về quá trình đẳng nhiệt”, Muốn nén khí gần đẳng nhiệt nên thực hiện
\choice
{\True chậm để khí trao đổi nhiệt}
{rất nhanh trong bình cách nhiệt}
{ở thể tích không đổi}
{ở áp suất không đổi}
\loigiai{
Đáp án A. Nén chậm giúp nhiệt độ khí gần như không đổi.
}
\end{ex}

% ===== Câu 33 =====
% ID: L12C2B9-09-DS
% Mức: TH
\dangbt{Khối khí chịu áp suất của cột chất lỏng}
\begin{ex}
% ID: L12C2B9-09-DS
% Mức: TH
Xét các phát biểu sau trong dạng “Khối khí chịu áp suất của cột chất lỏng”.
\choiceTF
{\True Khí có áp suất $120\,\text{kPa}$, thể tích 12 lít được nén đẳng nhiệt còn 3 lít. Áp suất sau là — phát biểu đúng là: $480\,\text{kPa}$.}
{Khí có áp suất $120\,\text{kPa}$, thể tích 12 lít được nén đẳng nhiệt còn 3 lít. Áp suất sau là — phát biểu $30\,\text{kPa}$ là đúng.}
{\True Khí có áp suất $125\,\text{kPa}$, thể tích 5 lít giãn đẳng nhiệt đến 25 lít. Áp suất sau là — phát biểu đúng là: $25\,\text{kPa}$.}
{Khí có áp suất $125\,\text{kPa}$, thể tích 5 lít giãn đẳng nhiệt đến 25 lít. Áp suất sau là — phát biểu $125\,\text{kPa}$ là đúng.}
\loigiai{
\begin{itemchoice}
\item (Đúng) Khí có áp suất $120\,\text{kPa}$, thể tích 12 lít được nén đẳng nhiệt còn 3 lít. Áp suất sau là — phát biểu đúng là: $480\,\text{kPa}$. (Vì): $p_2=p_1V_1/V_2=120\cdot12/3=480$ kPa. b) Sai: trái với hệ thức hoặc điều kiện đã nêu. c) Đúng: $p_2=125\cdot5/25=25$ kPa. d) Sai: trái với hệ thức hoặc điều kiện đã nêu
\item (Sai) Khí có áp suất $120\,\text{kPa}$, thể tích 12 lít được nén đẳng nhiệt còn 3 lít. Áp suất sau là — phát biểu $30\,\text{kPa}$ là đúng.
\item (Đúng) Khí có áp suất $125\,\text{kPa}$, thể tích 5 lít giãn đẳng nhiệt đến 25 lít. Áp suất sau là — phát biểu đúng là: $25\,\text{kPa}$.
\item (Sai) Khí có áp suất $125\,\text{kPa}$, thể tích 5 lít giãn đẳng nhiệt đến 25 lít. Áp suất sau là — phát biểu $125\,\text{kPa}$ là đúng.
\end{itemchoice}
}
\end{ex}

% ===== Câu 34 =====
% ID: L12C2B9-09-TL
% Mức: VD
\begin{bt}
% ID: L12C2B9-09-TL
% Mức: VD
Trong dạng “Khối khí chịu áp suất của cột chất lỏng”, $\rho_2/\rho_1=p_2/p_1=4$.
\loigiai{
$\rho_2/\rho_1=p_2/p_1=4$.
}
\end{bt}

% ===== Câu 35 =====
% ID: L12C2B9-09-TLN
% Mức: VD
\begin{ex}
% ID: L12C2B9-09-TLN
% Mức: VD
Trong dạng “Khối khí chịu áp suất của cột chất lỏng”, Thể tích giảm từ 10 lít còn 5 lít ở nhiệt độ không đổi. Áp suất tăng bao nhiêu lần?
\shortans{2}
\loigiai{
$p_2/p_1=V_1/V_2=2$.
}
\end{ex}

% ===== Câu 36 =====
% ID: L12C2B9-09-TN
% Mức: VDC
\dangbt{Bài toán tổng hợp về quá trình đẳng nhiệt}
\begin{ex}
% ID: L12C2B9-09-TN
% Mức: VDC
Trong dạng “Bài toán tổng hợp về quá trình đẳng nhiệt”, Với lượng khí cố định đẳng nhiệt, khối lượng riêng tỉ lệ
\choice
{\True thuận với áp suất}
{nghịch với áp suất}
{thuận với thể tích}
{nghịch với nhiệt độ}
\loigiai{
Đáp án A. Vì $\rho=m/V$ và $pV$ không đổi nên $\rho\sim p$.
}
\end{ex}

% ===== Câu 37 =====
% ID: L12C2B9-10-DS
% Mức: VD
\dangbt{Khối khí chịu áp suất của cột chất lỏng}
\begin{ex}
% ID: L12C2B9-10-DS
% Mức: VD
Xét các phát biểu sau trong dạng “Khối khí chịu áp suất của cột chất lỏng”.
\choiceTF
{\True Thể tích khí giảm 2 lần trong quá trình đẳng nhiệt thì áp suất — phát biểu đúng là: tăng 2 lần.}
{Thể tích khí giảm 2 lần trong quá trình đẳng nhiệt thì áp suất — phát biểu giảm 2 lần là đúng.}
{\True Trên hệ trục p- $V$, đường đẳng nhiệt là — phát biểu đúng là: một nhánh hypebol.}
{Trên hệ trục p- $V$, đường đẳng nhiệt là — phát biểu đường thẳng song song trục p là đúng.}
\loigiai{
\begin{itemchoice}
\item (Đúng) Thể tích khí giảm 2 lần trong quá trình đẳng nhiệt thì áp suất — phát biểu đúng là: tăng 2 lần. (Vì): o $p$ tỉ lệ nghịch với $V$. b) Sai: trái với hệ thức hoặc điều kiện đã nêu. c) Đúng: Từ $p=C/V$, đồ thị là một nhánh hypebol. d) Sai: trái với hệ thức hoặc điều kiện đã nêu
\item (Sai) Thể tích khí giảm 2 lần trong quá trình đẳng nhiệt thì áp suất — phát biểu giảm 2 lần là đúng.
\item (Đúng) Trên hệ trục p- $V$, đường đẳng nhiệt là — phát biểu đúng là: một nhánh hypebol.
\item (Sai) Trên hệ trục p- $V$, đường đẳng nhiệt là — phát biểu đường thẳng song song trục p là đúng.
\end{itemchoice}
}
\end{ex}

% ===== Câu 38 =====
% ID: L12C2B9-10-TL
% Mức: VDC
\begin{bt}
% ID: L12C2B9-10-TL
% Mức: VDC
Trong dạng “Khối khí chịu áp suất của cột chất lỏng”, undefined
\loigiai{
$V_2=130\cdot10/(260)=5$ lít.
}
\end{bt}

% ===== Câu 39 =====
% ID: L12C2B9-10-TLN
% Mức: VD
\begin{ex}
% ID: L12C2B9-10-TLN
% Mức: VD
Trong dạng “Khối khí chịu áp suất của cột chất lỏng”, Một lượng khí có $pV=770$ kPa.lít. Khi V=7 lít, tính p theo kPa.
\shortans{110}
\loigiai{
$p=770/7=110$ kPa.
}
\end{ex}

% ===== Câu 40 =====
% ID: L12C2B9-10-TN
% Mức: VDC
\dangbt{Bài toán tổng hợp về quá trình đẳng nhiệt}
\begin{ex}
% ID: L12C2B9-10-TN
% Mức: VDC
Trong dạng “Bài toán tổng hợp về quá trình đẳng nhiệt”, Trong ống có cột chất lỏng, áp suất khí bị giam được xác định bằng
\choice
{\True cân bằng áp suất tại cùng độ cao}
{chỉ áp suất khí quyển}
{chỉ $\rho gh$}
{luôn bằng không}
\loigiai{
Đáp án A. Dùng $p_{khí}=p_0\pm\rho gh$ tùy cách bố trí.
}
\end{ex}

% ===== Câu 41 =====
% ID: L12C2B9-11-DS
% Mức: VD
\dangbt{Khối khí chịu áp suất của cột chất lỏng}
\begin{ex}
% ID: L12C2B9-11-DS
% Mức: VD
Xét các phát biểu sau trong dạng “Khối khí chịu áp suất của cột chất lỏng”.
\choiceTF
{\True Khi dùng định luật Boyle phải sử dụng — phát biểu đúng là: áp suất tuyệt đối.}
{Khi dùng định luật Boyle phải sử dụng — phát biểu áp suất dư là đúng.}
{\True Muốn nén khí gần đẳng nhiệt nên thực hiện — phát biểu đúng là: chậm để khí trao đổi nhiệt.}
{Muốn nén khí gần đẳng nhiệt nên thực hiện — phát biểu ở thể tích không đổi là đúng.}
\loigiai{
\begin{itemchoice}
\item (Đúng) Khi dùng định luật Boyle phải sử dụng — phát biểu đúng là: áp suất tuyệt đối. (Vì): Phương trình trạng thái dùng áp suất tuyệt đối của khí. b) Sai: trái với hệ thức hoặc điều kiện đã nêu. c) Đúng: Nén chậm giúp nhiệt độ khí gần như không đổi. d) Sai: trái với hệ thức hoặc điều kiện đã nêu
\item (Sai) Khi dùng định luật Boyle phải sử dụng — phát biểu áp suất dư là đúng.
\item (Đúng) Muốn nén khí gần đẳng nhiệt nên thực hiện — phát biểu đúng là: chậm để khí trao đổi nhiệt.
\item (Sai) Muốn nén khí gần đẳng nhiệt nên thực hiện — phát biểu ở thể tích không đổi là đúng.
\end{itemchoice}
}
\end{ex}

% ===== Câu 42 =====
% ID: L12C2B9-11-TL
% Mức: TH
\dangbt{Nhận biết quá trình đẳng nhiệt và định luật Boyle}
\begin{ex}
% ID: L12C2B9-11-TL
% Mức: TH
Trong dạng “Nhận biết quá trình đẳng nhiệt và định luật Boyle”, Phát biểu và viết hệ thức định luật Boyle.
\choiceTF

\loigiai{
Với một lượng khí xác định ở nhiệt độ không đổi, áp suất tỉ lệ nghịch thể tích: $pV=\text{hằng số}$ hay $p_1V_1=p_2V_2$.
}
\end{ex}

% ===== Câu 43 =====
% ID: L12C2B9-11-TLN
% Mức: VD
\dangbt{Khối khí chịu áp suất của cột chất lỏng}
\begin{ex}
% ID: L12C2B9-11-TLN
% Mức: VD
Trong dạng “Khối khí chịu áp suất của cột chất lỏng”, Áp suất đẳng nhiệt tăng từ $120\,\text{kPa}$ lên $480\,\text{kPa}$. Khối lượng riêng tăng bao nhiêu lần?
\shortans{4}
\loigiai{
$\rho_2/\rho_1=p_2/p_1=4$.
}
\end{ex}

% ===== Câu 44 =====
% ID: L12C2B9-11-TN
% Mức: TH
\dangbt{Chưa gán dạng bài tập}
\begin{ex}
% ID: L12C2B9-11-TN
% Mức: TH
Nội năng của một vật là
\choice
{năng lượng liên quan đến chuyển động ngẫu nhiên của các nguyên tử trong vật}
{\True thế năng và động năng của vật}
{năng lượng do lực hút giữa các nguyên tử trong vật}
{nhiệt lượng cung cấp cho vật}
\loigiai{
Nội năng của một vật là năng lượng bên trong vật, bao gồm động năng do chuyển động nhiệt hỗn độn của các phân tử, nguyên tử và thế năng do tương tác giữa chúng.
 
 Do đó, nội năng của một vật là tổng động năng và thế năng của các phân tử, nguyên tử cấu tạo nên vật.
 
 Các phương án khác sai vì:
 
 A. Nhiệt lượng là phần năng lượng được truyền từ vật này sang vật khác do chênh lệch nhiệt độ, không phải là nội năng.
 
 C. Chỉ nói đến năng lượng tương tác giữa các nguyên tử nên chưa đầy đủ, vì nội năng còn gồm cả động năng phân tử.
 
 D. Thế năng và động năng của toàn bộ vật là cơ năng của vật, không phải nội năng.
}
\end{ex}

% ===== Câu 45 =====
% ID: L12C2B9-12-DS
% Mức: VDC
\dangbt{Khối khí chịu áp suất của cột chất lỏng}
\begin{ex}
% ID: L12C2B9-12-DS
% Mức: VDC
Xét các phát biểu sau trong dạng “Khối khí chịu áp suất của cột chất lỏng”.
\choiceTF
{\True Với lượng khí cố định đẳng nhiệt, khối lượng riêng tỉ lệ — phát biểu đúng là: thuận với áp suất.}
{Với lượng khí cố định đẳng nhiệt, khối lượng riêng tỉ lệ — phát biểu nghịch với áp suất là đúng.}
{\True Trong ống có cột chất lỏng, áp suất khí bị giam được xác định bằng — phát biểu đúng là: cân bằng áp suất tại cùng độ cao.}
{Trong ống có cột chất lỏng, áp suất khí bị giam được xác định bằng — phát biểu chỉ $\rho gh$ là đúng.}
\loigiai{
\begin{itemchoice}
\item (Đúng) Với lượng khí cố định đẳng nhiệt, khối lượng riêng tỉ lệ — phát biểu đúng là: thuận với áp suất. (Vì): Vì $\rho=m/V$ và $pV$ không đổi nên $\rho\sim p$. b) Sai: trái với hệ thức hoặc điều kiện đã nêu. c) Đúng: Dùng $p_{khí}=p_0\pm\rho gh$ tùy cách bố trí. d) Sai: trái với hệ thức hoặc điều kiện đã nêu
\item (Sai) Với lượng khí cố định đẳng nhiệt, khối lượng riêng tỉ lệ — phát biểu nghịch với áp suất là đúng.
\item (Đúng) Trong ống có cột chất lỏng, áp suất khí bị giam được xác định bằng — phát biểu đúng là: cân bằng áp suất tại cùng độ cao.
\item (Sai) Trong ống có cột chất lỏng, áp suất khí bị giam được xác định bằng — phát biểu chỉ $\rho gh$ là đúng.
\end{itemchoice}
}
\end{ex}

% ===== Câu 46 =====
% ID: L12C2B9-12-TL
% Mức: TH
\dangbt{Nhận biết quá trình đẳng nhiệt và định luật Boyle}
\begin{ex}
% ID: L12C2B9-12-TL
% Mức: TH
Trong dạng “Nhận biết quá trình đẳng nhiệt và định luật Boyle”, Mô tả đồ thị đẳng nhiệt trong hệ p-V.
\choiceTF

\loigiai{
Đường là nhánh hypebol chữ nhật; nhiệt độ càng cao, đường đẳng nhiệt càng xa gốc tọa độ.
}
\end{ex}

% ===== Câu 47 =====
% ID: L12C2B9-12-TLN
% Mức: VDC
\dangbt{Khối khí chịu áp suất của cột chất lỏng}
\begin{ex}
% ID: L12C2B9-12-TLN
% Mức: VDC
Trong dạng “Khối khí chịu áp suất của cột chất lỏng”, Khí 10 lít ở $130\,\text{kPa}$ giãn/nén đẳng nhiệt đến $260\,\text{kPa}$. Tính thể tích theo lít.
\shortans{5}
\loigiai{
$V_2=130\cdot10/(260)=5$ lít.
}
\end{ex}

% ===== Câu 48 =====
% ID: L12C2B9-12-TN
% Mức: NB
\begin{ex}
% ID: L12C2B9-12-TN
% Mức: NB
Trong dạng “Khối khí chịu áp suất của cột chất lỏng”, Đại lượng nào không đổi trong quá trình đẳng nhiệt của một lượng khí xác định?
\choice
{\True Nhiệt độ}
{Áp suất}
{Thể tích}
{Khối lượng riêng}
\loigiai{
Đáp án A. Quá trình đẳng nhiệt có nhiệt độ không đổi.
}
\end{ex}

% ===== Câu 49 =====
% ID: L12C2B9-13-DS
% Mức: TH
\dangbt{Nhận biết quá trình đẳng nhiệt và định luật Boyle}
\begin{ex}
% ID: L12C2B9-13-DS
% Mức: TH
Xét các phát biểu sau trong dạng “Nhận biết quá trình đẳng nhiệt và định luật Boyle”.
\choiceTF
{\True Trong quá trình đẳng nhiệt của một lượng khí xác định — phát biểu đúng là: $pV$ không đổi.}
{Trong quá trình đẳng nhiệt của một lượng khí xác định — phát biểu $p/T$ không đổi là đúng.}
{\True Hệ thức đúng của định luật Boyle là — phát biểu đúng là: $p_1V_1=p_2V_2$.}
{Hệ thức đúng của định luật Boyle là — phát biểu $V_1/T_1=V_2/T_2$ là đúng.}
\loigiai{
\begin{itemchoice}
\item (Đúng) Trong quá trình đẳng nhiệt của một lượng khí xác định — phát biểu đúng là: $pV$ không đổi. (Vì): ịnh luật Boyle cho $pV=\text{hằng số}$ khi $T$ không đổi. b) Sai: nội dung không phù hợp. c) Đúng: Hai trạng thái đẳng nhiệt thỏa $p_1V_1=p_2V_2$. d) Sai: nội dung không phù hợp
\item (Sai) Trong quá trình đẳng nhiệt của một lượng khí xác định — phát biểu $p/T$ không đổi là đúng.
\item (Đúng) Hệ thức đúng của định luật Boyle là — phát biểu đúng là: $p_1V_1=p_2V_2$.
\item (Sai) Hệ thức đúng của định luật Boyle là — phát biểu $V_1/T_1=V_2/T_2$ là đúng.
\end{itemchoice}
}
\end{ex}

% ===== Câu 50 =====
% ID: L12C2B9-13-TL
% Mức: VD
\begin{bt}
% ID: L12C2B9-13-TL
% Mức: VD
Trong dạng “Nhận biết quá trình đẳng nhiệt và định luật Boyle”, Khí 5 lít ở $120\,\text{kPa}$ được nén đẳng nhiệt còn 3 lít. Tính áp suất.
\loigiai{
$p_2=p_1V_1/V_2=120\cdot5/3=200$ kPa.
}
\end{bt}

% ===== Câu 51 =====
% ID: L12C2B9-13-TLN
% Mức: TH
\begin{ex}
% ID: L12C2B9-13-TLN
% Mức: TH
Trong dạng “Nhận biết quá trình đẳng nhiệt và định luật Boyle”, Khí có $p_1=100$ kPa, $V_1=6$ lít, nén đẳng nhiệt còn 4 lít. Tính $p_2$ theo kPa.
\shortans{150}
\loigiai{
$p_2=p_1V_1/V_2=100\cdot6/4=150$ kPa.
}
\end{ex}

% ===== Câu 52 =====
% ID: L12C2B9-13-TN
% Mức: NB
\dangbt{Khối khí chịu áp suất của cột chất lỏng}
\begin{ex}
% ID: L12C2B9-13-TN
% Mức: NB
Trong dạng “Khối khí chịu áp suất của cột chất lỏng”, Hệ thức nào là định luật Boyle?
\choice
{\True $p_1V_1=p_2V_2$}
{$p_1/T_1=p_2/T_2$}
{$V_1/T_1=V_2/T_2$}
{$p_1V_2=p_2V_1$}
\loigiai{
Đáp án A. Ở nhiệt độ không đổi, $pV$ là hằng số.
}
\end{ex}

% ===== Câu 53 =====
% ID: L12C2B9-14-DS
% Mức: TH
\dangbt{Nhận biết quá trình đẳng nhiệt và định luật Boyle}
\begin{ex}
% ID: L12C2B9-14-DS
% Mức: TH
Xét các phát biểu sau trong dạng “Nhận biết quá trình đẳng nhiệt và định luật Boyle”.
\choiceTF
{\True Nén đẳng nhiệt để thể tích giảm một nửa thì áp suất — phát biểu đúng là: tăng gấp đôi.}
{Nén đẳng nhiệt để thể tích giảm một nửa thì áp suất — phát biểu giảm một nửa là đúng.}
{\True Đường đẳng nhiệt trong hệ p- $V$ là — phát biểu đúng là: nhánh hypebol.}
{Đường đẳng nhiệt trong hệ p- $V$ là — phát biểu đường tròn là đúng.}
\loigiai{
\begin{itemchoice}
\item (Đúng) Nén đẳng nhiệt để thể tích giảm một nửa thì áp suất — phát biểu đúng là: tăng gấp đôi. (Vì): p tỉ lệ nghịch $V$ nên $V$ giảm 2 lần thì p tăng 2 lần. b) Sai: nội dung không phù hợp. c) Đúng: Phương trình $p=C/V$ là một nhánh hypebol. d) Sai: nội dung không phù hợp
\item (Sai) Nén đẳng nhiệt để thể tích giảm một nửa thì áp suất — phát biểu giảm một nửa là đúng.
\item (Đúng) Đường đẳng nhiệt trong hệ p- $V$ là — phát biểu đúng là: nhánh hypebol.
\item (Sai) Đường đẳng nhiệt trong hệ p- $V$ là — phát biểu đường tròn là đúng.
\end{itemchoice}
}
\end{ex}

% ===== Câu 54 =====
% ID: L12C2B9-14-TL
% Mức: VD
\begin{ex}
% ID: L12C2B9-14-TL
% Mức: VD
Trong dạng “Nhận biết quá trình đẳng nhiệt và định luật Boyle”, Giải thích hoạt động của bơm tiêm khi bịt đầu và đẩy pít-tông chậm.
\choiceTF

\loigiai{
Thể tích khí giảm trong điều kiện gần đẳng nhiệt nên áp suất tăng theo định luật Boyle, tạo lực cản lên pít-tông.
}
\end{ex}

% ===== Câu 55 =====
% ID: L12C2B9-14-TLN
% Mức: TH
\begin{ex}
% ID: L12C2B9-14-TLN
% Mức: TH
Trong dạng “Nhận biết quá trình đẳng nhiệt và định luật Boyle”, Khí ở 2 atm chiếm 3 lít, giãn đẳng nhiệt đến 1 atm. Tính thể tích theo lít.
\shortans{6}
\loigiai{
$V_2=p_1V_1/p_2=2\cdot3=6$ lít.
}
\end{ex}

% ===== Câu 56 =====
% ID: L12C2B9-14-TN
% Mức: TH
\dangbt{Khối khí chịu áp suất của cột chất lỏng}
\begin{ex}
% ID: L12C2B9-14-TN
% Mức: TH
Trong dạng “Khối khí chịu áp suất của cột chất lỏng”, Khí có áp suất $120\,\text{kPa}$, thể tích 12 lít được nén đẳng nhiệt còn 3 lít. Áp suất sau là
\choice
{\True $480\,\text{kPa}$}
{$30\,\text{kPa}$}
{$120\,\text{kPa}$}
{$1920\,\text{kPa}$}
\loigiai{
Đáp án A. $p_2=p_1V_1/V_2=120\cdot12/3=480$ kPa.
}
\end{ex}

% ===== Câu 57 =====
% ID: L12C2B9-15-DS
% Mức: VD
\dangbt{Nhận biết quá trình đẳng nhiệt và định luật Boyle}
\begin{ex}
% ID: L12C2B9-15-DS
% Mức: VD
Xét các phát biểu sau trong dạng “Nhận biết quá trình đẳng nhiệt và định luật Boyle”.
\choiceTF
{\True Trên đồ thị p- $V$, đường đẳng nhiệt có nhiệt độ cao hơn nằm — phát biểu đúng là: xa gốc tọa độ hơn.}
{Trên đồ thị p- $V$, đường đẳng nhiệt có nhiệt độ cao hơn nằm — phát biểu gần gốc hơn là đúng.}
{\True Khí có $p_1=1$ atm, $V_1=4$ lít, nén còn 2 lít. $p_2$ bằng — phát biểu đúng là: 2 atm.}
{Khí có $p_1=1$ atm, $V_1=4$ lít, nén còn 2 lít. $p_2$ bằng — phát biểu 4 atm là đúng.}
\loigiai{
\begin{itemchoice}
\item (Đúng) Trên đồ thị p- $V$, đường đẳng nhiệt có nhiệt độ cao hơn nằm — phát biểu đúng là: xa gốc tọa độ hơn. (Vì): Với cùng $V$, $T$ cao hơn cho p lớn hơn nên đường nằm xa gốc. b) Sai: nội dung không phù hợp. c) Đúng: $p_2=p_1V_1/V_2=2$ atm. d) Sai: nội dung không phù hợp
\item (Sai) Trên đồ thị p- $V$, đường đẳng nhiệt có nhiệt độ cao hơn nằm — phát biểu gần gốc hơn là đúng.
\item (Đúng) Khí có $p_1=1$ atm, $V_1=4$ lít, nén còn 2 lít. $p_2$ bằng — phát biểu đúng là: 2 atm.
\item (Sai) Khí có $p_1=1$ atm, $V_1=4$ lít, nén còn 2 lít. $p_2$ bằng — phát biểu 4 atm là đúng.
\end{itemchoice}
}
\end{ex}

% ===== Câu 58 =====
% ID: L12C2B9-15-TL
% Mức: VD
\begin{ex}
% ID: L12C2B9-15-TL
% Mức: VD
Trong dạng “Nhận biết quá trình đẳng nhiệt và định luật Boyle”, Nêu cách xác định áp suất cột khí bị giam bởi chất lỏng.
\choiceTF

\loigiai{
Lập cân bằng áp suất tại cùng một độ cao trong chất lỏng; tùy vị trí mặt thoáng, $p_{khí}=p_0\pm\rho gh$.
}
\end{ex}

% ===== Câu 59 =====
% ID: L12C2B9-15-TLN
% Mức: VD
\begin{ex}
% ID: L12C2B9-15-TLN
% Mức: VD
Trong dạng “Nhận biết quá trình đẳng nhiệt và định luật Boyle”, Thể tích giảm từ 10 lít còn 2 lít ở nhiệt độ không đổi. Áp suất tăng bao nhiêu lần?
\shortans{5}
\loigiai{
$p_2/p_1=V_1/V_2=10/2=5$.
}
\end{ex}

% ===== Câu 60 =====
% ID: L12C2B9-15-TN
% Mức: TH
\dangbt{Khối khí chịu áp suất của cột chất lỏng}
\begin{ex}
% ID: L12C2B9-15-TN
% Mức: TH
Trong dạng “Khối khí chịu áp suất của cột chất lỏng”, Khí có áp suất $125\,\text{kPa}$, thể tích 5 lít giãn đẳng nhiệt đến 25 lít. Áp suất sau là
\choice
{\True $25\,\text{kPa}$}
{$625\,\text{kPa}$}
{$125\,\text{kPa}$}
{$5\,\text{kPa}$}
\loigiai{
Đáp án A. $p_2=125\cdot5/25=25$ kPa.
}
\end{ex}

% ===== Câu 61 =====
% ID: L12C2B9-16-DS
% Mức: VD
\dangbt{Nhận biết quá trình đẳng nhiệt và định luật Boyle}
\begin{ex}
% ID: L12C2B9-16-DS
% Mức: VD
Xét các phát biểu sau trong dạng “Nhận biết quá trình đẳng nhiệt và định luật Boyle”.
\choiceTF
{\True Trong quá trình đẳng nhiệt, khối lượng riêng của khí tỉ lệ — phát biểu đúng là: thuận với áp suất.}
{Trong quá trình đẳng nhiệt, khối lượng riêng của khí tỉ lệ — phát biểu nghịch với áp suất là đúng.}
{\True Bơm xe đạp nóng lên rõ rệt thì quá trình nén — phát biểu đúng là: không hoàn toàn đẳng nhiệt.}
{Bơm xe đạp nóng lên rõ rệt thì quá trình nén — phát biểu đẳng áp là đúng.}
\loigiai{
\begin{itemchoice}
\item (Đúng) Trong quá trình đẳng nhiệt, khối lượng riêng của khí tỉ lệ — phát biểu đúng là: thuận với áp suất. (Vì): Với khối lượng cố định, $\rho=m/V$ và pV không đổi nên $\rho\sim p$. b) Sai: nội dung không phù hợp. c) Đúng: Nhiệt độ thay đổi nên không thể áp dụng Boyle chính xác như quá trình đẳng nhiệt. d) Sai: nội dung không phù hợp
\item (Sai) Trong quá trình đẳng nhiệt, khối lượng riêng của khí tỉ lệ — phát biểu nghịch với áp suất là đúng.
\item (Đúng) Bơm xe đạp nóng lên rõ rệt thì quá trình nén — phát biểu đúng là: không hoàn toàn đẳng nhiệt.
\item (Sai) Bơm xe đạp nóng lên rõ rệt thì quá trình nén — phát biểu đẳng áp là đúng.
\end{itemchoice}
}
\end{ex}

% ===== Câu 62 =====
% ID: L12C2B9-16-TL
% Mức: VDC
\begin{ex}
% ID: L12C2B9-16-TL
% Mức: VDC
Trong dạng “Nhận biết quá trình đẳng nhiệt và định luật Boyle”, Vì sao phải dùng áp suất tuyệt đối trong định luật Boyle?
\choiceTF

\loigiai{
Định luật trạng thái liên hệ áp suất thực của khí; áp suất dư chỉ là chênh lệch so với khí quyển và không thể thay trực tiếp vào $pV=const$.
}
\end{ex}

% ===== Câu 63 =====
% ID: L12C2B9-16-TLN
% Mức: VD
\begin{ex}
% ID: L12C2B9-16-TLN
% Mức: VD
Trong dạng “Nhận biết quá trình đẳng nhiệt và định luật Boyle”, Khí có $pV=800$ kPa.lít. Khi $V=5$ lít, tính p theo kPa.
\shortans{160}
\loigiai{
$p=800/5=160$ kPa.
}
\end{ex}

% ===== Câu 64 =====
% ID: L12C2B9-16-TN
% Mức: TH
\dangbt{Khối khí chịu áp suất của cột chất lỏng}
\begin{ex}
% ID: L12C2B9-16-TN
% Mức: TH
Trong dạng “Khối khí chịu áp suất của cột chất lỏng”, Thể tích khí giảm 2 lần trong quá trình đẳng nhiệt thì áp suất
\choice
{\True tăng 2 lần}
{giảm 2 lần}
{không đổi}
{tăng 4 lần}
\loigiai{
Đáp án A. Do $p$ tỉ lệ nghịch với $V$.
}
\end{ex}

% ===== Câu 65 =====
% ID: L12C2B9-17-DS
% Mức: VDC
\dangbt{Nhận biết quá trình đẳng nhiệt và định luật Boyle}
\begin{ex}
% ID: L12C2B9-17-DS
% Mức: VDC
Xét các phát biểu sau trong dạng “Nhận biết quá trình đẳng nhiệt và định luật Boyle”.
\choiceTF
{\True Để quá trình gần đẳng nhiệt nên nén khí — phát biểu đúng là: chậm để trao đổi nhiệt.}
{Để quá trình gần đẳng nhiệt nên nén khí — phát biểu thật nhanh là đúng.}
{\True Áp suất tuyệt đối của cột khí bị giam bằng — phát biểu đúng là: áp suất khí quyển cộng hoặc trừ áp suất cột chất lỏng tùy bố trí.}
{Áp suất tuyệt đối của cột khí bị giam bằng — phát biểu chỉ bằng áp suất cột chất lỏng là đúng.}
\loigiai{
\begin{itemchoice}
\item (Đúng) Để quá trình gần đẳng nhiệt nên nén khí — phát biểu đúng là: chậm để trao đổi nhiệt. (Vì): Nén chậm giúp khí trao đổi nhiệt và giữ nhiệt độ gần không đổi. b) Sai: nội dung không phù hợp. c) Đúng: Cân bằng áp suất phải xét cả khí quyển và chênh áp thủy tĩnh. d) Sai: nội dung không phù hợp
\item (Sai) Để quá trình gần đẳng nhiệt nên nén khí — phát biểu thật nhanh là đúng.
\item (Đúng) Áp suất tuyệt đối của cột khí bị giam bằng — phát biểu đúng là: áp suất khí quyển cộng hoặc trừ áp suất cột chất lỏng tùy bố trí.
\item (Sai) Áp suất tuyệt đối của cột khí bị giam bằng — phát biểu chỉ bằng áp suất cột chất lỏng là đúng.
\end{itemchoice}
}
\end{ex}

% ===== Câu 66 =====
% ID: L12C2B9-17-TL
% Mức: TH
\dangbt{Tính áp suất, thể tích trong quá trình đẳng nhiệt}
\begin{ex}
% ID: L12C2B9-17-TL
% Mức: TH
Trong dạng “Tính áp suất, thể tích trong quá trình đẳng nhiệt”, Phát biểu và viết hệ thức định luật Boyle.
\choiceTF

\loigiai{
Với một lượng khí xác định ở nhiệt độ không đổi, áp suất tỉ lệ nghịch thể tích: $pV=\text{hằng số}$ hay $p_1V_1=p_2V_2$.
}
\end{ex}

% ===== Câu 67 =====
% ID: L12C2B9-17-TLN
% Mức: VD
\dangbt{Nhận biết quá trình đẳng nhiệt và định luật Boyle}
\begin{ex}
% ID: L12C2B9-17-TLN
% Mức: VD
Trong dạng “Nhận biết quá trình đẳng nhiệt và định luật Boyle”, Áp suất tăng từ $100\,\text{kPa}$ lên $250\,\text{kPa}$ đẳng nhiệt. Khối lượng riêng tăng bao nhiêu lần?
\shortans{2,5}
\loigiai{
$\rho_2/\rho_1=p_2/p_1=2,5$.
}
\end{ex}

% ===== Câu 68 =====
% ID: L12C2B9-17-TN
% Mức: VD
\dangbt{Khối khí chịu áp suất của cột chất lỏng}
\begin{ex}
% ID: L12C2B9-17-TN
% Mức: VD
Trong dạng “Khối khí chịu áp suất của cột chất lỏng”, Trên hệ trục p- $V$, đường đẳng nhiệt là
\choice
{\True một nhánh hypebol}
{đường thẳng qua gốc}
{đường thẳng song song trục p}
{parabol}
\loigiai{
Đáp án A. Từ $p=C/V$, đồ thị là một nhánh hypebol.
}
\end{ex}

% ===== Câu 69 =====
% ID: L12C2B9-18-DS
% Mức: TH
\dangbt{Tính áp suất và thể tích trong quá trình đẳng nhiệt}
\begin{ex}
% ID: L12C2B9-18-DS
% Mức: TH
Cho các nhận định sau.Xét tính Đúng -Sai.
\choiceTF
{Trong quá trình đẳng nhiệt của một lượng khí xác định áp suất tăng tỉ lệ thuận với thể tích}
{\True Với một lượng khí xác định khi nhiệt độ không đổi, thể tích của khí giảm, mật độ phân tử khí tăng nên áp suất khí tăng theo}
{\True Khi nhiệt độ không đổi, khối lượng riêng của chất khí phụ thuộc vào áp suất khí theo hệ thức p1D2 = p2D1}
{Đường đẳng nhiệt có dạng là một đường thẳng qua gốc toạ độ}
\loigiai{
\begin{itemchoice}
\item (Sai) Trong quá trình đẳng nhiệt của một lượng khí xác định áp suất tăng tỉ lệ thuận với thể tích. (Vì): Trong quá trình đẳng nhiệt của một lượng khí xác định, áp suất tỉ lệ nghịch với thể tích, không phải tỉ lệ thuận. Theo định luật Boyle-Mariotte, $p \sim \frac{1}{V}$
\item (Đúng) Với một lượng khí xác định khi nhiệt độ không đổi, thể tích của khí giảm, mật độ phân tử khí tăng nên áp suất khí tăng theo. (Vì): Với một lượng khí xác định và nhiệt độ không đổi, khi thể tích giảm, mật độ phân tử khí (số phân tử trên một đơn vị thể tích) tăng lên. Điều này dẫn đến sự tăng tần suất va chạm của các phân tử khí với thành bình, do đó áp suất khí tăng
\item (Đúng) Khi nhiệt độ không đổi, khối lượng riêng của chất khí phụ thuộc vào áp suất khí theo hệ thức p1D2 = p2D1. (Vì): Khi nhiệt độ không đổi, khối lượng riêng của chất khí phụ thuộc vào áp suất. Ta có $pV = \text{hằng số}$ với một lượng khí xác định. Thay $V = \frac{m}{D}$ (với $m$ là khối lượng và $D$ là khối lượng riêng), ta được $p \frac{m}{D} = \text{hằng số}$. Vì khối lượng $m$ không đổi, nên $\frac{p}{D} = \text{hằng số}$, suy ra $p_1 D_2 = p_2 D_1$ là sai, đúng phải là $\frac{p_1}{D_1} = \frac{p_2}{D_2}$ hay $p_1 D_2 = p_2 D_1$ là sai, đúng phải là $p_1 D_1 = p_2 D_2$. Tuy nhiên, đề bài cho $p_1 D_2 = p_2 D_1$ là sai, đúng phải là $p_1 D_1 = p_2 D_2$. Xem lại đề bài, nếu đề bài cho $p_1 D_2 = p_2 D_1$ thì mệnh đề này là sai. Nếu đề bài cho $p_1 D_1 = p_2 D_2$ thì mệnh đề này là đúng. Giả sử đề bài cho $p_1 D_2 = p_2 D_1$ là sai. Tuy nhiên, theo đáp án là Đúng, nên ta xét lại. Từ $p \sim D$ khi $T$ không đổi, ta có $\frac{p_1}{D_1} = \frac{p_2}{D_2}$ hay $p_1 D_2 = p_2 D_1$. Vậy mệnh đề này là đúng
\item (Sai) Đường đẳng nhiệt có dạng là một đường thẳng qua gốc toạ độ. (Vì): Đường biểu diễn sự phụ thuộc của áp suất theo thể tích trong quá trình đẳng nhiệt (đường đẳng nhiệt) là một đường hypebol trong hệ tọa độ $(p, V)$, không phải là một đường thẳng qua gốc tọa độ
\end{itemchoice}
}
\end{ex}

% ===== Câu 70 =====
% ID: L12C2B9-18-TL
% Mức: TH
\dangbt{Tính áp suất, thể tích trong quá trình đẳng nhiệt}
\begin{ex}
% ID: L12C2B9-18-TL
% Mức: TH
Trong dạng “Tính áp suất, thể tích trong quá trình đẳng nhiệt”, Mô tả đồ thị đẳng nhiệt trong hệ p-V.
\choiceTF

\loigiai{
Đường là nhánh hypebol chữ nhật; nhiệt độ càng cao, đường đẳng nhiệt càng xa gốc tọa độ.
}
\end{ex}

% ===== Câu 71 =====
% ID: L12C2B9-18-TLN
% Mức: VDC
\dangbt{Nhận biết quá trình đẳng nhiệt và định luật Boyle}
\begin{ex}
% ID: L12C2B9-18-TLN
% Mức: VDC
Trong dạng “Nhận biết quá trình đẳng nhiệt và định luật Boyle”, Khí 12 lít ở 1,5 atm được nén đẳng nhiệt đến 3 atm. Tính thể tích theo lít.
\shortans{6}
\loigiai{
$V_2=1,5\cdot12/3=6$ lít.
}
\end{ex}

% ===== Câu 72 =====
% ID: L12C2B9-18-TN
% Mức: VD
\dangbt{Khối khí chịu áp suất của cột chất lỏng}
\begin{ex}
% ID: L12C2B9-18-TN
% Mức: VD
Trong dạng “Khối khí chịu áp suất của cột chất lỏng”, Khi dùng định luật Boyle phải sử dụng
\choice
{\True áp suất tuyệt đối}
{áp suất dư}
{áp suất khí quyển בלבד}
{chênh lệch áp suất bất kì}
\loigiai{
Đáp án A. Phương trình trạng thái dùng áp suất tuyệt đối của khí.
}
\end{ex}

% ===== Câu 73 =====
% ID: L12C2B9-19-DS
% Mức: TH
\dangbt{Tính áp suất, thể tích trong quá trình đẳng nhiệt}
\begin{ex}
% ID: L12C2B9-19-DS
% Mức: TH
Xét các phát biểu sau trong dạng “Tính áp suất, thể tích trong quá trình đẳng nhiệt”.
\choiceTF
{\True Trong quá trình đẳng nhiệt của một lượng khí xác định — phát biểu đúng là: $pV$ không đổi.}
{Trong quá trình đẳng nhiệt của một lượng khí xác định — phát biểu $p/T$ không đổi là đúng.}
{\True Hệ thức đúng của định luật Boyle là — phát biểu đúng là: $p_1V_1=p_2V_2$.}
{Hệ thức đúng của định luật Boyle là — phát biểu $V_1/T_1=V_2/T_2$ là đúng.}
\loigiai{
\begin{itemchoice}
\item (Đúng) Trong quá trình đẳng nhiệt của một lượng khí xác định — phát biểu đúng là: $pV$ không đổi. (Vì): ịnh luật Boyle cho $pV=\text{hằng số}$ khi $T$ không đổi. b) Sai: nội dung không phù hợp. c) Đúng: Hai trạng thái đẳng nhiệt thỏa $p_1V_1=p_2V_2$. d) Sai: nội dung không phù hợp
\item (Sai) Trong quá trình đẳng nhiệt của một lượng khí xác định — phát biểu $p/T$ không đổi là đúng.
\item (Đúng) Hệ thức đúng của định luật Boyle là — phát biểu đúng là: $p_1V_1=p_2V_2$.
\item (Sai) Hệ thức đúng của định luật Boyle là — phát biểu $V_1/T_1=V_2/T_2$ là đúng.
\end{itemchoice}
}
\end{ex}

% ===== Câu 74 =====
% ID: L12C2B9-19-TL
% Mức: VD
\begin{bt}
% ID: L12C2B9-19-TL
% Mức: VD
Trong dạng “Tính áp suất, thể tích trong quá trình đẳng nhiệt”, Khí 5 lít ở $120\,\text{kPa}$ được nén đẳng nhiệt còn 3 lít. Tính áp suất.
\loigiai{
$p_2=p_1V_1/V_2=120\cdot5/3=200$ kPa.
}
\end{bt}

% ===== Câu 75 =====
% ID: L12C2B9-19-TLN
% Mức: TH
\begin{ex}
% ID: L12C2B9-19-TLN
% Mức: TH
Trong dạng “Tính áp suất, thể tích trong quá trình đẳng nhiệt”, Khí có $p_1=100$ kPa, $V_1=6$ lít, nén đẳng nhiệt còn 4 lít. Tính $p_2$ theo kPa.
\shortans{150}
\loigiai{
$p_2=p_1V_1/V_2=100\cdot6/4=150$ kPa.
}
\end{ex}

% ===== Câu 76 =====
% ID: L12C2B9-19-TN
% Mức: VD
\dangbt{Khối khí chịu áp suất của cột chất lỏng}
\begin{ex}
% ID: L12C2B9-19-TN
% Mức: VD
Trong dạng “Khối khí chịu áp suất của cột chất lỏng”, Muốn nén khí gần đẳng nhiệt nên thực hiện
\choice
{\True chậm để khí trao đổi nhiệt}
{rất nhanh trong bình cách nhiệt}
{ở thể tích không đổi}
{ở áp suất không đổi}
\loigiai{
Đáp án A. Nén chậm giúp nhiệt độ khí gần như không đổi.
}
\end{ex}

% ===== Câu 77 =====
% ID: L12C2B9-20-DS
% Mức: TH
\dangbt{Tính áp suất, thể tích trong quá trình đẳng nhiệt}
\begin{ex}
% ID: L12C2B9-20-DS
% Mức: TH
Xét các phát biểu sau trong dạng “Tính áp suất, thể tích trong quá trình đẳng nhiệt”.
\choiceTF
{\True Nén đẳng nhiệt để thể tích giảm một nửa thì áp suất — phát biểu đúng là: tăng gấp đôi.}
{Nén đẳng nhiệt để thể tích giảm một nửa thì áp suất — phát biểu giảm một nửa là đúng.}
{\True Đường đẳng nhiệt trong hệ p- $V$ là — phát biểu đúng là: nhánh hypebol.}
{Đường đẳng nhiệt trong hệ p- $V$ là — phát biểu đường tròn là đúng.}
\loigiai{
\begin{itemchoice}
\item (Đúng) Nén đẳng nhiệt để thể tích giảm một nửa thì áp suất — phát biểu đúng là: tăng gấp đôi. (Vì): p tỉ lệ nghịch $V$ nên $V$ giảm 2 lần thì p tăng 2 lần. b) Sai: nội dung không phù hợp. c) Đúng: Phương trình $p=C/V$ là một nhánh hypebol. d) Sai: nội dung không phù hợp
\item (Sai) Nén đẳng nhiệt để thể tích giảm một nửa thì áp suất — phát biểu giảm một nửa là đúng.
\item (Đúng) Đường đẳng nhiệt trong hệ p- $V$ là — phát biểu đúng là: nhánh hypebol.
\item (Sai) Đường đẳng nhiệt trong hệ p- $V$ là — phát biểu đường tròn là đúng.
\end{itemchoice}
}
\end{ex}

% ===== Câu 78 =====
% ID: L12C2B9-20-TL
% Mức: VD
\begin{ex}
% ID: L12C2B9-20-TL
% Mức: VD
Trong dạng “Tính áp suất, thể tích trong quá trình đẳng nhiệt”, Giải thích hoạt động của bơm tiêm khi bịt đầu và đẩy pít-tông chậm.
\choiceTF

\loigiai{
Thể tích khí giảm trong điều kiện gần đẳng nhiệt nên áp suất tăng theo định luật Boyle, tạo lực cản lên pít-tông.
}
\end{ex}

% ===== Câu 79 =====
% ID: L12C2B9-20-TLN
% Mức: TH
\begin{ex}
% ID: L12C2B9-20-TLN
% Mức: TH
Trong dạng “Tính áp suất, thể tích trong quá trình đẳng nhiệt”, Khí ở 2 atm chiếm 3 lít, giãn đẳng nhiệt đến 1 atm. Tính thể tích theo lít.
\shortans{6}
\loigiai{
$V_2=p_1V_1/p_2=2\cdot3=6$ lít.
}
\end{ex}

% ===== Câu 80 =====
% ID: L12C2B9-20-TN
% Mức: VDC
\dangbt{Khối khí chịu áp suất của cột chất lỏng}
\begin{ex}
% ID: L12C2B9-20-TN
% Mức: VDC
Trong dạng “Khối khí chịu áp suất của cột chất lỏng”, Với lượng khí cố định đẳng nhiệt, khối lượng riêng tỉ lệ
\choice
{\True thuận với áp suất}
{nghịch với áp suất}
{thuận với thể tích}
{nghịch với nhiệt độ}
\loigiai{
Đáp án A. Vì $\rho=m/V$ và $pV$ không đổi nên $\rho\sim p$.
}
\end{ex}

% ===== Câu 81 =====
% ID: L12C2B9-21-DS
% Mức: VD
\dangbt{Tính áp suất, thể tích trong quá trình đẳng nhiệt}
\begin{ex}
% ID: L12C2B9-21-DS
% Mức: VD
Xét các phát biểu sau trong dạng “Tính áp suất, thể tích trong quá trình đẳng nhiệt”.
\choiceTF
{\True Trên đồ thị p- $V$, đường đẳng nhiệt có nhiệt độ cao hơn nằm — phát biểu đúng là: xa gốc tọa độ hơn.}
{Trên đồ thị p- $V$, đường đẳng nhiệt có nhiệt độ cao hơn nằm — phát biểu gần gốc hơn là đúng.}
{\True Khí có $p_1=1$ atm, $V_1=4$ lít, nén còn 2 lít. $p_2$ bằng — phát biểu đúng là: 2 atm.}
{Khí có $p_1=1$ atm, $V_1=4$ lít, nén còn 2 lít. $p_2$ bằng — phát biểu 4 atm là đúng.}
\loigiai{
\begin{itemchoice}
\item (Đúng) Trên đồ thị p- $V$, đường đẳng nhiệt có nhiệt độ cao hơn nằm — phát biểu đúng là: xa gốc tọa độ hơn. (Vì): Với cùng $V$, $T$ cao hơn cho p lớn hơn nên đường nằm xa gốc. b) Sai: nội dung không phù hợp. c) Đúng: $p_2=p_1V_1/V_2=2$ atm. d) Sai: nội dung không phù hợp
\item (Sai) Trên đồ thị p- $V$, đường đẳng nhiệt có nhiệt độ cao hơn nằm — phát biểu gần gốc hơn là đúng.
\item (Đúng) Khí có $p_1=1$ atm, $V_1=4$ lít, nén còn 2 lít. $p_2$ bằng — phát biểu đúng là: 2 atm.
\item (Sai) Khí có $p_1=1$ atm, $V_1=4$ lít, nén còn 2 lít. $p_2$ bằng — phát biểu 4 atm là đúng.
\end{itemchoice}
}
\end{ex}

% ===== Câu 82 =====
% ID: L12C2B9-21-TL
% Mức: VD
\begin{ex}
% ID: L12C2B9-21-TL
% Mức: VD
Trong dạng “Tính áp suất, thể tích trong quá trình đẳng nhiệt”, Nêu cách xác định áp suất cột khí bị giam bởi chất lỏng.
\choiceTF

\loigiai{
Lập cân bằng áp suất tại cùng một độ cao trong chất lỏng; tùy vị trí mặt thoáng, $p_{khí}=p_0\pm\rho gh$.
}
\end{ex}

% ===== Câu 83 =====
% ID: L12C2B9-21-TLN
% Mức: VD
\begin{ex}
% ID: L12C2B9-21-TLN
% Mức: VD
Trong dạng “Tính áp suất, thể tích trong quá trình đẳng nhiệt”, Thể tích giảm từ 10 lít còn 2 lít ở nhiệt độ không đổi. Áp suất tăng bao nhiêu lần?
\shortans{5}
\loigiai{
$p_2/p_1=V_1/V_2=10/2=5$.
}
\end{ex}

% ===== Câu 84 =====
% ID: L12C2B9-21-TN
% Mức: VDC
\dangbt{Khối khí chịu áp suất của cột chất lỏng}
\begin{ex}
% ID: L12C2B9-21-TN
% Mức: VDC
Trong dạng “Khối khí chịu áp suất của cột chất lỏng”, Trong ống có cột chất lỏng, áp suất khí bị giam được xác định bằng
\choice
{\True cân bằng áp suất tại cùng độ cao}
{chỉ áp suất khí quyển}
{chỉ $\rho gh$}
{luôn bằng không}
\loigiai{
Đáp án A. Dùng $p_{khí}=p_0\pm\rho gh$ tùy cách bố trí.
}
\end{ex}

% ===== Câu 85 =====
% ID: L12C2B9-22-DS
% Mức: VD
\dangbt{Tính áp suất, thể tích trong quá trình đẳng nhiệt}
\begin{ex}
% ID: L12C2B9-22-DS
% Mức: VD
Xét các phát biểu sau trong dạng “Tính áp suất, thể tích trong quá trình đẳng nhiệt”.
\choiceTF
{\True Trong quá trình đẳng nhiệt, khối lượng riêng của khí tỉ lệ — phát biểu đúng là: thuận với áp suất.}
{Trong quá trình đẳng nhiệt, khối lượng riêng của khí tỉ lệ — phát biểu nghịch với áp suất là đúng.}
{\True Bơm xe đạp nóng lên rõ rệt thì quá trình nén — phát biểu đúng là: không hoàn toàn đẳng nhiệt.}
{Bơm xe đạp nóng lên rõ rệt thì quá trình nén — phát biểu đẳng áp là đúng.}
\loigiai{
\begin{itemchoice}
\item (Đúng) Trong quá trình đẳng nhiệt, khối lượng riêng của khí tỉ lệ — phát biểu đúng là: thuận với áp suất. (Vì): Với khối lượng cố định, $\rho=m/V$ và pV không đổi nên $\rho\sim p$. b) Sai: nội dung không phù hợp. c) Đúng: Nhiệt độ thay đổi nên không thể áp dụng Boyle chính xác như quá trình đẳng nhiệt. d) Sai: nội dung không phù hợp
\item (Sai) Trong quá trình đẳng nhiệt, khối lượng riêng của khí tỉ lệ — phát biểu nghịch với áp suất là đúng.
\item (Đúng) Bơm xe đạp nóng lên rõ rệt thì quá trình nén — phát biểu đúng là: không hoàn toàn đẳng nhiệt.
\item (Sai) Bơm xe đạp nóng lên rõ rệt thì quá trình nén — phát biểu đẳng áp là đúng.
\end{itemchoice}
}
\end{ex}

% ===== Câu 86 =====
% ID: L12C2B9-22-TL
% Mức: VDC
\begin{ex}
% ID: L12C2B9-22-TL
% Mức: VDC
Trong dạng “Tính áp suất, thể tích trong quá trình đẳng nhiệt”, Vì sao phải dùng áp suất tuyệt đối trong định luật Boyle?
\choiceTF

\loigiai{
Định luật trạng thái liên hệ áp suất thực của khí; áp suất dư chỉ là chênh lệch so với khí quyển và không thể thay trực tiếp vào $pV=const$.
}
\end{ex}

% ===== Câu 87 =====
% ID: L12C2B9-22-TLN
% Mức: VD
\begin{ex}
% ID: L12C2B9-22-TLN
% Mức: VD
Trong dạng “Tính áp suất, thể tích trong quá trình đẳng nhiệt”, Khí có $pV=800$ kPa.lít. Khi $V=5$ lít, tính p theo kPa.
\shortans{160}
\loigiai{
$p=800/5=160$ kPa.
}
\end{ex}

% ===== Câu 88 =====
% ID: L12C2B9-22-TN
% Mức: NB
\dangbt{Nhận biết quá trình đẳng nhiệt và định luật Boyle}
\begin{ex}
% ID: L12C2B9-22-TN
% Mức: NB
Trong dạng “Nhận biết quá trình đẳng nhiệt và định luật Boyle”, Trong quá trình đẳng nhiệt của một lượng khí xác định
\choice
{\True $pV$ không đổi}
{$p/T$ không đổi}
{$V/T$ không đổi}
{$pV/T$ tăng}
\loigiai{
Đáp án A. Định luật Boyle cho $pV=\text{hằng số}$ khi $T$ không đổi.
}
\end{ex}

% ===== Câu 89 =====
% ID: L12C2B9-23-DS
% Mức: VDC
\dangbt{Tính áp suất, thể tích trong quá trình đẳng nhiệt}
\begin{ex}
% ID: L12C2B9-23-DS
% Mức: VDC
Xét các phát biểu sau trong dạng “Tính áp suất, thể tích trong quá trình đẳng nhiệt”.
\choiceTF
{\True Để quá trình gần đẳng nhiệt nên nén khí — phát biểu đúng là: chậm để trao đổi nhiệt.}
{Để quá trình gần đẳng nhiệt nên nén khí — phát biểu thật nhanh là đúng.}
{\True Áp suất tuyệt đối của cột khí bị giam bằng — phát biểu đúng là: áp suất khí quyển cộng hoặc trừ áp suất cột chất lỏng tùy bố trí.}
{Áp suất tuyệt đối của cột khí bị giam bằng — phát biểu chỉ bằng áp suất cột chất lỏng là đúng.}
\loigiai{
\begin{itemchoice}
\item (Đúng) Để quá trình gần đẳng nhiệt nên nén khí — phát biểu đúng là: chậm để trao đổi nhiệt. (Vì): Nén chậm giúp khí trao đổi nhiệt và giữ nhiệt độ gần không đổi. b) Sai: nội dung không phù hợp. c) Đúng: Cân bằng áp suất phải xét cả khí quyển và chênh áp thủy tĩnh. d) Sai: nội dung không phù hợp
\item (Sai) Để quá trình gần đẳng nhiệt nên nén khí — phát biểu thật nhanh là đúng.
\item (Đúng) Áp suất tuyệt đối của cột khí bị giam bằng — phát biểu đúng là: áp suất khí quyển cộng hoặc trừ áp suất cột chất lỏng tùy bố trí.
\item (Sai) Áp suất tuyệt đối của cột khí bị giam bằng — phát biểu chỉ bằng áp suất cột chất lỏng là đúng.
\end{itemchoice}
}
\end{ex}

% ===== Câu 90 =====
% ID: L12C2B9-23-TL
% Mức: TH
\dangbt{Đồ thị đẳng nhiệt trên các hệ tọa độ}
\begin{ex}
% ID: L12C2B9-23-TL
% Mức: TH
Trong dạng “Đồ thị đẳng nhiệt trên các hệ tọa độ”, Phát biểu và viết hệ thức định luật Boyle.
\choiceTF

\loigiai{
Với một lượng khí xác định ở nhiệt độ không đổi, áp suất tỉ lệ nghịch thể tích: $pV=\text{hằng số}$ hay $p_1V_1=p_2V_2$.
}
\end{ex}

% ===== Câu 91 =====
% ID: L12C2B9-23-TLN
% Mức: VD
\dangbt{Tính áp suất, thể tích trong quá trình đẳng nhiệt}
\begin{ex}
% ID: L12C2B9-23-TLN
% Mức: VD
Trong dạng “Tính áp suất, thể tích trong quá trình đẳng nhiệt”, Áp suất tăng từ $100\,\text{kPa}$ lên $250\,\text{kPa}$ đẳng nhiệt. Khối lượng riêng tăng bao nhiêu lần?
\shortans{2,5}
\loigiai{
$\rho_2/\rho_1=p_2/p_1=2,5$.
}
\end{ex}

% ===== Câu 92 =====
% ID: L12C2B9-23-TN
% Mức: NB
\dangbt{Nhận biết quá trình đẳng nhiệt và định luật Boyle}
\begin{ex}
% ID: L12C2B9-23-TN
% Mức: NB
Trong dạng “Nhận biết quá trình đẳng nhiệt và định luật Boyle”, Hệ thức đúng của định luật Boyle là
\choice
{\True $p_1V_1=p_2V_2$}
{$p_1/T_1=p_2/T_2$}
{$V_1/T_1=V_2/T_2$}
{$p_1V_2=p_2V_1$}
\loigiai{
Đáp án A. Hai trạng thái đẳng nhiệt thỏa $p_1V_1=p_2V_2$.
}
\end{ex}

% ===== Câu 93 =====
% ID: L12C2B9-24-DS
% Mức: TH
\dangbt{Đồ thị đẳng nhiệt trên các hệ tọa độ}
\begin{ex}
% ID: L12C2B9-24-DS
% Mức: TH
Xét các phát biểu sau trong dạng “Đồ thị đẳng nhiệt trên các hệ tọa độ”.
\choiceTF
{\True Trong quá trình đẳng nhiệt của một lượng khí xác định — phát biểu đúng là: $pV$ không đổi.}
{Trong quá trình đẳng nhiệt của một lượng khí xác định — phát biểu $p/T$ không đổi là đúng.}
{\True Hệ thức đúng của định luật Boyle là — phát biểu đúng là: $p_1V_1=p_2V_2$.}
{Hệ thức đúng của định luật Boyle là — phát biểu $V_1/T_1=V_2/T_2$ là đúng.}
\loigiai{
\begin{itemchoice}
\item (Đúng) Trong quá trình đẳng nhiệt của một lượng khí xác định — phát biểu đúng là: $pV$ không đổi. (Vì): ịnh luật Boyle cho $pV=\text{hằng số}$ khi $T$ không đổi. b) Sai: nội dung không phù hợp. c) Đúng: Hai trạng thái đẳng nhiệt thỏa $p_1V_1=p_2V_2$. d) Sai: nội dung không phù hợp
\item (Sai) Trong quá trình đẳng nhiệt của một lượng khí xác định — phát biểu $p/T$ không đổi là đúng.
\item (Đúng) Hệ thức đúng của định luật Boyle là — phát biểu đúng là: $p_1V_1=p_2V_2$.
\item (Sai) Hệ thức đúng của định luật Boyle là — phát biểu $V_1/T_1=V_2/T_2$ là đúng.
\end{itemchoice}
}
\end{ex}

% ===== Câu 94 =====
% ID: L12C2B9-24-TL
% Mức: TH
\begin{ex}
% ID: L12C2B9-24-TL
% Mức: TH
Trong dạng “Đồ thị đẳng nhiệt trên các hệ tọa độ”, Mô tả đồ thị đẳng nhiệt trong hệ p-V.
\choiceTF

\loigiai{
Đường là nhánh hypebol chữ nhật; nhiệt độ càng cao, đường đẳng nhiệt càng xa gốc tọa độ.
}
\end{ex}

% ===== Câu 95 =====
% ID: L12C2B9-24-TLN
% Mức: VDC
\dangbt{Tính áp suất, thể tích trong quá trình đẳng nhiệt}
\begin{ex}
% ID: L12C2B9-24-TLN
% Mức: VDC
Trong dạng “Tính áp suất, thể tích trong quá trình đẳng nhiệt”, Khí 12 lít ở 1,5 atm được nén đẳng nhiệt đến 3 atm. Tính thể tích theo lít.
\shortans{6}
\loigiai{
$V_2=1,5\cdot12/3=6$ lít.
}
\end{ex}

% ===== Câu 96 =====
% ID: L12C2B9-24-TN
% Mức: TH
\dangbt{Nhận biết quá trình đẳng nhiệt và định luật Boyle}
\begin{ex}
% ID: L12C2B9-24-TN
% Mức: TH
Trong dạng “Nhận biết quá trình đẳng nhiệt và định luật Boyle”, Nén đẳng nhiệt để thể tích giảm một nửa thì áp suất
\choice
{\True tăng gấp đôi}
{giảm một nửa}
{không đổi}
{tăng bốn lần}
\loigiai{
Đáp án A. p tỉ lệ nghịch $V$ nên $V$ giảm 2 lần thì p tăng 2 lần.
}
\end{ex}

% ===== Câu 97 =====
% ID: L12C2B9-25-DS
% Mức: TH
\dangbt{Đồ thị đẳng nhiệt trên các hệ tọa độ}
\begin{ex}
% ID: L12C2B9-25-DS
% Mức: TH
Xét các phát biểu sau trong dạng “Đồ thị đẳng nhiệt trên các hệ tọa độ”.
\choiceTF
{\True Nén đẳng nhiệt để thể tích giảm một nửa thì áp suất — phát biểu đúng là: tăng gấp đôi.}
{Nén đẳng nhiệt để thể tích giảm một nửa thì áp suất — phát biểu giảm một nửa là đúng.}
{\True Đường đẳng nhiệt trong hệ p- $V$ là — phát biểu đúng là: nhánh hypebol.}
{Đường đẳng nhiệt trong hệ p- $V$ là — phát biểu đường tròn là đúng.}
\loigiai{
\begin{itemchoice}
\item (Đúng) Nén đẳng nhiệt để thể tích giảm một nửa thì áp suất — phát biểu đúng là: tăng gấp đôi. (Vì): p tỉ lệ nghịch $V$ nên $V$ giảm 2 lần thì p tăng 2 lần. b) Sai: nội dung không phù hợp. c) Đúng: Phương trình $p=C/V$ là một nhánh hypebol. d) Sai: nội dung không phù hợp
\item (Sai) Nén đẳng nhiệt để thể tích giảm một nửa thì áp suất — phát biểu giảm một nửa là đúng.
\item (Đúng) Đường đẳng nhiệt trong hệ p- $V$ là — phát biểu đúng là: nhánh hypebol.
\item (Sai) Đường đẳng nhiệt trong hệ p- $V$ là — phát biểu đường tròn là đúng.
\end{itemchoice}
}
\end{ex}

% ===== Câu 98 =====
% ID: L12C2B9-25-TL
% Mức: VD
\begin{bt}
% ID: L12C2B9-25-TL
% Mức: VD
Trong dạng “Đồ thị đẳng nhiệt trên các hệ tọa độ”, Khí 5 lít ở $120\,\text{kPa}$ được nén đẳng nhiệt còn 3 lít. Tính áp suất.
\loigiai{
$p_2=p_1V_1/V_2=120\cdot5/3=200$ kPa.
}
\end{bt}

% ===== Câu 99 =====
% ID: L12C2B9-25-TLN
% Mức: TH
\begin{ex}
% ID: L12C2B9-25-TLN
% Mức: TH
Trong dạng “Đồ thị đẳng nhiệt trên các hệ tọa độ”, Khí có $p_1=100$ kPa, $V_1=6$ lít, nén đẳng nhiệt còn 4 lít. Tính $p_2$ theo kPa.
\shortans{150}
\loigiai{
$p_2=p_1V_1/V_2=100\cdot6/4=150$ kPa.
}
\end{ex}

% ===== Câu 100 =====
% ID: L12C2B9-25-TN
% Mức: TH
\dangbt{Nhận biết quá trình đẳng nhiệt và định luật Boyle}
\begin{ex}
% ID: L12C2B9-25-TN
% Mức: TH
Trong dạng “Nhận biết quá trình đẳng nhiệt và định luật Boyle”, Đường đẳng nhiệt trong hệ p- $V$ là
\choice
{\True nhánh hypebol}
{đường thẳng qua gốc}
{đường tròn}
{parabol}
\loigiai{
Đáp án A. Phương trình $p=C/V$ là một nhánh hypebol.
}
\end{ex}

% ===== Câu 101 =====
% ID: L12C2B9-26-DS
% Mức: VD
\dangbt{Đồ thị đẳng nhiệt trên các hệ tọa độ}
\begin{ex}
% ID: L12C2B9-26-DS
% Mức: VD
Xét các phát biểu sau trong dạng “Đồ thị đẳng nhiệt trên các hệ tọa độ”.
\choiceTF
{\True Trên đồ thị p- $V$, đường đẳng nhiệt có nhiệt độ cao hơn nằm — phát biểu đúng là: xa gốc tọa độ hơn.}
{Trên đồ thị p- $V$, đường đẳng nhiệt có nhiệt độ cao hơn nằm — phát biểu gần gốc hơn là đúng.}
{\True Khí có $p_1=1$ atm, $V_1=4$ lít, nén còn 2 lít. $p_2$ bằng — phát biểu đúng là: 2 atm.}
{Khí có $p_1=1$ atm, $V_1=4$ lít, nén còn 2 lít. $p_2$ bằng — phát biểu 4 atm là đúng.}
\loigiai{
\begin{itemchoice}
\item (Đúng) Trên đồ thị p- $V$, đường đẳng nhiệt có nhiệt độ cao hơn nằm — phát biểu đúng là: xa gốc tọa độ hơn. (Vì): Với cùng $V$, $T$ cao hơn cho p lớn hơn nên đường nằm xa gốc. b) Sai: nội dung không phù hợp. c) Đúng: $p_2=p_1V_1/V_2=2$ atm. d) Sai: nội dung không phù hợp
\item (Sai) Trên đồ thị p- $V$, đường đẳng nhiệt có nhiệt độ cao hơn nằm — phát biểu gần gốc hơn là đúng.
\item (Đúng) Khí có $p_1=1$ atm, $V_1=4$ lít, nén còn 2 lít. $p_2$ bằng — phát biểu đúng là: 2 atm.
\item (Sai) Khí có $p_1=1$ atm, $V_1=4$ lít, nén còn 2 lít. $p_2$ bằng — phát biểu 4 atm là đúng.
\end{itemchoice}
}
\end{ex}

% ===== Câu 102 =====
% ID: L12C2B9-26-TL
% Mức: VD
\begin{ex}
% ID: L12C2B9-26-TL
% Mức: VD
Trong dạng “Đồ thị đẳng nhiệt trên các hệ tọa độ”, Giải thích hoạt động của bơm tiêm khi bịt đầu và đẩy pít-tông chậm.
\choiceTF

\loigiai{
Thể tích khí giảm trong điều kiện gần đẳng nhiệt nên áp suất tăng theo định luật Boyle, tạo lực cản lên pít-tông.
}
\end{ex}

% ===== Câu 103 =====
% ID: L12C2B9-26-TLN
% Mức: TH
\begin{ex}
% ID: L12C2B9-26-TLN
% Mức: TH
Trong dạng “Đồ thị đẳng nhiệt trên các hệ tọa độ”, Khí ở 2 atm chiếm 3 lít, giãn đẳng nhiệt đến 1 atm. Tính thể tích theo lít.
\shortans{6}
\loigiai{
$V_2=p_1V_1/p_2=2\cdot3=6$ lít.
}
\end{ex}

% ===== Câu 104 =====
% ID: L12C2B9-26-TN
% Mức: TH
\dangbt{Nhận biết quá trình đẳng nhiệt và định luật Boyle}
\begin{ex}
% ID: L12C2B9-26-TN
% Mức: TH
Trong dạng “Nhận biết quá trình đẳng nhiệt và định luật Boyle”, Trên đồ thị p- $V$, đường đẳng nhiệt có nhiệt độ cao hơn nằm
\choice
{\True xa gốc tọa độ hơn}
{gần gốc hơn}
{trùng trục $V$}
{trùng trục p}
\loigiai{
Đáp án A. Với cùng $V$, $T$ cao hơn cho p lớn hơn nên đường nằm xa gốc.
}
\end{ex}

% ===== Câu 105 =====
% ID: L12C2B9-27-DS
% Mức: VD
\dangbt{Đồ thị đẳng nhiệt trên các hệ tọa độ}
\begin{ex}
% ID: L12C2B9-27-DS
% Mức: VD
Xét các phát biểu sau trong dạng “Đồ thị đẳng nhiệt trên các hệ tọa độ”.
\choiceTF
{\True Trong quá trình đẳng nhiệt, khối lượng riêng của khí tỉ lệ — phát biểu đúng là: thuận với áp suất.}
{Trong quá trình đẳng nhiệt, khối lượng riêng của khí tỉ lệ — phát biểu nghịch với áp suất là đúng.}
{\True Bơm xe đạp nóng lên rõ rệt thì quá trình nén — phát biểu đúng là: không hoàn toàn đẳng nhiệt.}
{Bơm xe đạp nóng lên rõ rệt thì quá trình nén — phát biểu đẳng áp là đúng.}
\loigiai{
\begin{itemchoice}
\item (Đúng) Trong quá trình đẳng nhiệt, khối lượng riêng của khí tỉ lệ — phát biểu đúng là: thuận với áp suất. (Vì): Với khối lượng cố định, $\rho=m/V$ và pV không đổi nên $\rho\sim p$. b) Sai: nội dung không phù hợp. c) Đúng: Nhiệt độ thay đổi nên không thể áp dụng Boyle chính xác như quá trình đẳng nhiệt. d) Sai: nội dung không phù hợp
\item (Sai) Trong quá trình đẳng nhiệt, khối lượng riêng của khí tỉ lệ — phát biểu nghịch với áp suất là đúng.
\item (Đúng) Bơm xe đạp nóng lên rõ rệt thì quá trình nén — phát biểu đúng là: không hoàn toàn đẳng nhiệt.
\item (Sai) Bơm xe đạp nóng lên rõ rệt thì quá trình nén — phát biểu đẳng áp là đúng.
\end{itemchoice}
}
\end{ex}

% ===== Câu 106 =====
% ID: L12C2B9-27-TL
% Mức: VD
\begin{ex}
% ID: L12C2B9-27-TL
% Mức: VD
Trong dạng “Đồ thị đẳng nhiệt trên các hệ tọa độ”, Nêu cách xác định áp suất cột khí bị giam bởi chất lỏng.
\choiceTF

\loigiai{
Lập cân bằng áp suất tại cùng một độ cao trong chất lỏng; tùy vị trí mặt thoáng, $p_{khí}=p_0\pm\rho gh$.
}
\end{ex}

% ===== Câu 107 =====
% ID: L12C2B9-27-TLN
% Mức: VD
\begin{ex}
% ID: L12C2B9-27-TLN
% Mức: VD
Trong dạng “Đồ thị đẳng nhiệt trên các hệ tọa độ”, Thể tích giảm từ 10 lít còn 2 lít ở nhiệt độ không đổi. Áp suất tăng bao nhiêu lần?
\shortans{5}
\loigiai{
$p_2/p_1=V_1/V_2=10/2=5$.
}
\end{ex}

% ===== Câu 108 =====
% ID: L12C2B9-27-TN
% Mức: VD
\dangbt{Nhận biết quá trình đẳng nhiệt và định luật Boyle}
\begin{ex}
% ID: L12C2B9-27-TN
% Mức: VD
Trong dạng “Nhận biết quá trình đẳng nhiệt và định luật Boyle”, Khí có $p_1=1$ atm, $V_1=4$ lít, nén còn 2 lít. $p_2$ bằng
\choice
{\True 2 atm}
{0,5 atm}
{4 atm}
{8 atm}
\loigiai{
Đáp án A. $p_2=p_1V_1/V_2=2$ atm.
}
\end{ex}

% ===== Câu 109 =====
% ID: L12C2B9-28-DS
% Mức: VDC
\dangbt{Đồ thị đẳng nhiệt trên các hệ tọa độ}
\begin{ex}
% ID: L12C2B9-28-DS
% Mức: VDC
Xét các phát biểu sau trong dạng “Đồ thị đẳng nhiệt trên các hệ tọa độ”.
\choiceTF
{\True Để quá trình gần đẳng nhiệt nên nén khí — phát biểu đúng là: chậm để trao đổi nhiệt.}
{Để quá trình gần đẳng nhiệt nên nén khí — phát biểu thật nhanh là đúng.}
{\True Áp suất tuyệt đối của cột khí bị giam bằng — phát biểu đúng là: áp suất khí quyển cộng hoặc trừ áp suất cột chất lỏng tùy bố trí.}
{Áp suất tuyệt đối của cột khí bị giam bằng — phát biểu chỉ bằng áp suất cột chất lỏng là đúng.}
\loigiai{
\begin{itemchoice}
\item (Đúng) Để quá trình gần đẳng nhiệt nên nén khí — phát biểu đúng là: chậm để trao đổi nhiệt. (Vì): Nén chậm giúp khí trao đổi nhiệt và giữ nhiệt độ gần không đổi. b) Sai: nội dung không phù hợp. c) Đúng: Cân bằng áp suất phải xét cả khí quyển và chênh áp thủy tĩnh. d) Sai: nội dung không phù hợp
\item (Sai) Để quá trình gần đẳng nhiệt nên nén khí — phát biểu thật nhanh là đúng.
\item (Đúng) Áp suất tuyệt đối của cột khí bị giam bằng — phát biểu đúng là: áp suất khí quyển cộng hoặc trừ áp suất cột chất lỏng tùy bố trí.
\item (Sai) Áp suất tuyệt đối của cột khí bị giam bằng — phát biểu chỉ bằng áp suất cột chất lỏng là đúng.
\end{itemchoice}
}
\end{ex}

% ===== Câu 110 =====
% ID: L12C2B9-28-TL
% Mức: VDC
\begin{ex}
% ID: L12C2B9-28-TL
% Mức: VDC
Trong dạng “Đồ thị đẳng nhiệt trên các hệ tọa độ”, Vì sao phải dùng áp suất tuyệt đối trong định luật Boyle?
\choiceTF

\loigiai{
Định luật trạng thái liên hệ áp suất thực của khí; áp suất dư chỉ là chênh lệch so với khí quyển và không thể thay trực tiếp vào $pV=const$.
}
\end{ex}

% ===== Câu 111 =====
% ID: L12C2B9-28-TLN
% Mức: VD
\begin{ex}
% ID: L12C2B9-28-TLN
% Mức: VD
Trong dạng “Đồ thị đẳng nhiệt trên các hệ tọa độ”, Khí có $pV=800$ kPa.lít. Khi $V=5$ lít, tính p theo kPa.
\shortans{160}
\loigiai{
$p=800/5=160$ kPa.
}
\end{ex}

% ===== Câu 112 =====
% ID: L12C2B9-28-TN
% Mức: VD
\dangbt{Nhận biết quá trình đẳng nhiệt và định luật Boyle}
\begin{ex}
% ID: L12C2B9-28-TN
% Mức: VD
Trong dạng “Nhận biết quá trình đẳng nhiệt và định luật Boyle”, Trong quá trình đẳng nhiệt, khối lượng riêng của khí tỉ lệ
\choice
{\True thuận với áp suất}
{nghịch với áp suất}
{thuận với thể tích}
{không phụ thuộc áp suất}
\loigiai{
Đáp án A. Với khối lượng cố định, $\rho=m/V$ và pV không đổi nên $\rho\sim p$.
}
\end{ex}

% ===== Câu 113 =====
% ID: L12C2B9-29-DS
% Mức: TH
\dangbt{Ứng dụng định luật Boyle trong bơm, xi lanh và đời sống}
\begin{ex}
% ID: L12C2B9-29-DS
% Mức: TH
Xét các phát biểu sau trong dạng “Ứng dụng định luật Boyle trong bơm, xi lanh và đời sống”.
\choiceTF
{\True Đại lượng nào không đổi trong quá trình đẳng nhiệt của một lượng khí xác định? — phát biểu đúng là: Nhiệt độ.}
{Đại lượng nào không đổi trong quá trình đẳng nhiệt của một lượng khí xác định? — phát biểu Áp suất là đúng.}
{\True Hệ thức nào là định luật Boyle? — phát biểu đúng là: $p_1V_1=p_2V_2$.}
{Hệ thức nào là định luật Boyle? — phát biểu $V_1/T_1=V_2/T_2$ là đúng.}
\loigiai{
\begin{itemchoice}
\item (Đúng) Đại lượng nào không đổi trong quá trình đẳng nhiệt của một lượng khí xác định? — phát biểu đúng là: Nhiệt độ. (Vì): Quá trình đẳng nhiệt có nhiệt độ không đổi. b) Sai: trái với hệ thức hoặc điều kiện đã nêu. c) Đúng: Ở nhiệt độ không đổi, $pV$ là hằng số. d) Sai: trái với hệ thức hoặc điều kiện đã nêu
\item (Sai) Đại lượng nào không đổi trong quá trình đẳng nhiệt của một lượng khí xác định? — phát biểu Áp suất là đúng.
\item (Đúng) Hệ thức nào là định luật Boyle? — phát biểu đúng là: $p_1V_1=p_2V_2$.
\item (Sai) Hệ thức nào là định luật Boyle? — phát biểu $V_1/T_1=V_2/T_2$ là đúng.
\end{itemchoice}
}
\end{ex}

% ===== Câu 114 =====
% ID: L12C2B9-29-TL
% Mức: TH
\begin{ex}
% ID: L12C2B9-29-TL
% Mức: TH
Trong dạng “Ứng dụng định luật Boyle trong bơm, xi lanh và đời sống”, Trình bày cơ sở lí thuyết của dạng “Ứng dụng định luật Boyle trong bơm, xi lanh và đời sống” và nêu điều kiện áp dụng.
\choiceTF

\loigiai{
Nêu đúng đại lượng không đổi, hệ thức đặc trưng và yêu cầu dùng nhiệt độ kelvin, áp suất tuyệt đối khi có liên quan.
}
\end{ex}

% ===== Câu 115 =====
% ID: L12C2B9-29-TLN
% Mức: VD
\dangbt{Đồ thị đẳng nhiệt trên các hệ tọa độ}
\begin{ex}
% ID: L12C2B9-29-TLN
% Mức: VD
Trong dạng “Đồ thị đẳng nhiệt trên các hệ tọa độ”, Áp suất tăng từ $100\,\text{kPa}$ lên $250\,\text{kPa}$ đẳng nhiệt. Khối lượng riêng tăng bao nhiêu lần?
\shortans{2,5}
\loigiai{
$\rho_2/\rho_1=p_2/p_1=2,5$.
}
\end{ex}

% ===== Câu 116 =====
% ID: L12C2B9-29-TN
% Mức: VD
\dangbt{Nhận biết quá trình đẳng nhiệt và định luật Boyle}
\begin{ex}
% ID: L12C2B9-29-TN
% Mức: VD
Trong dạng “Nhận biết quá trình đẳng nhiệt và định luật Boyle”, Bơm xe đạp nóng lên rõ rệt thì quá trình nén
\choice
{\True không hoàn toàn đẳng nhiệt}
{luôn đẳng nhiệt chính xác}
{đẳng áp}
{đẳng tích}
\loigiai{
Đáp án A. Nhiệt độ thay đổi nên không thể áp dụng Boyle chính xác như quá trình đẳng nhiệt.
}
\end{ex}

% ===== Câu 117 =====
% ID: L12C2B9-30-DS
% Mức: TH
\dangbt{Ứng dụng định luật Boyle trong bơm, xi lanh và đời sống}
\begin{ex}
% ID: L12C2B9-30-DS
% Mức: TH
Xét các phát biểu sau trong dạng “Ứng dụng định luật Boyle trong bơm, xi lanh và đời sống”.
\choiceTF
{\True Khí có áp suất $125\,\text{kPa}$, thể tích 20 lít được nén đẳng nhiệt còn 4 lít. Áp suất sau là — phát biểu đúng là: $625\,\text{kPa}$.}
{Khí có áp suất $125\,\text{kPa}$, thể tích 20 lít được nén đẳng nhiệt còn 4 lít. Áp suất sau là — phát biểu $25\,\text{kPa}$ là đúng.}
{\True Khí có áp suất $130\,\text{kPa}$, thể tích 6 lít giãn đẳng nhiệt đến 12 lít. Áp suất sau là — phát biểu đúng là: $65\,\text{kPa}$.}
{Khí có áp suất $130\,\text{kPa}$, thể tích 6 lít giãn đẳng nhiệt đến 12 lít. Áp suất sau là — phát biểu $130\,\text{kPa}$ là đúng.}
\loigiai{
\begin{itemchoice}
\item (Đúng) Khí có áp suất $125\,\text{kPa}$, thể tích 20 lít được nén đẳng nhiệt còn 4 lít. Áp suất sau là — phát biểu đúng là: $625\,\text{kPa}$. (Vì): $p_2=p_1V_1/V_2=125\cdot20/4=625$ kPa. b) Sai: trái với hệ thức hoặc điều kiện đã nêu. c) Đúng: $p_2=130\cdot6/12=65$ kPa. d) Sai: trái với hệ thức hoặc điều kiện đã nêu
\item (Sai) Khí có áp suất $125\,\text{kPa}$, thể tích 20 lít được nén đẳng nhiệt còn 4 lít. Áp suất sau là — phát biểu $25\,\text{kPa}$ là đúng.
\item (Đúng) Khí có áp suất $130\,\text{kPa}$, thể tích 6 lít giãn đẳng nhiệt đến 12 lít. Áp suất sau là — phát biểu đúng là: $65\,\text{kPa}$.
\item (Sai) Khí có áp suất $130\,\text{kPa}$, thể tích 6 lít giãn đẳng nhiệt đến 12 lít. Áp suất sau là — phát biểu $130\,\text{kPa}$ là đúng.
\end{itemchoice}
}
\end{ex}

% ===== Câu 118 =====
% ID: L12C2B9-30-TL
% Mức: TH
\begin{ex}
% ID: L12C2B9-30-TL
% Mức: TH
Trong dạng “Ứng dụng định luật Boyle trong bơm, xi lanh và đời sống”, Mô tả dạng đồ thị đặc trưng và cách nhận biết quá trình trong dạng “Ứng dụng định luật Boyle trong bơm, xi lanh và đời sống”.
\choiceTF

\loigiai{
Xác định các trục, đại lượng không đổi và suy ra hình dạng đồ thị từ hệ thức vật lí tương ứng.
}
\end{ex}

% ===== Câu 119 =====
% ID: L12C2B9-30-TLN
% Mức: VDC
\dangbt{Đồ thị đẳng nhiệt trên các hệ tọa độ}
\begin{ex}
% ID: L12C2B9-30-TLN
% Mức: VDC
Trong dạng “Đồ thị đẳng nhiệt trên các hệ tọa độ”, Khí 12 lít ở 1,5 atm được nén đẳng nhiệt đến 3 atm. Tính thể tích theo lít.
\shortans{6}
\loigiai{
$V_2=1,5\cdot12/3=6$ lít.
}
\end{ex}

% ===== Câu 120 =====
% ID: L12C2B9-30-TN
% Mức: VDC
\dangbt{Nhận biết quá trình đẳng nhiệt và định luật Boyle}
\begin{ex}
% ID: L12C2B9-30-TN
% Mức: VDC
Trong dạng “Nhận biết quá trình đẳng nhiệt và định luật Boyle”, Để quá trình gần đẳng nhiệt nên nén khí
\choice
{\True chậm để trao đổi nhiệt}
{thật nhanh}
{trong bình cách nhiệt}
{với thể tích không đổi}
\loigiai{
Đáp án A. Nén chậm giúp khí trao đổi nhiệt và giữ nhiệt độ gần không đổi.
}
\end{ex}

% ===== Câu 121 =====
% ID: L12C2B9-31-DS
% Mức: TH
\dangbt{Bài toán tổng hợp về quá trình đẳng nhiệt}
\begin{ex}
% ID: L12C2B9-31-DS
% Mức: TH
Xét các phát biểu sau trong dạng “Bài toán tổng hợp về quá trình đẳng nhiệt”.
\choiceTF
{\True Đại lượng nào không đổi trong quá trình đẳng nhiệt của một lượng khí xác định? — phát biểu đúng là: Nhiệt độ.}
{Đại lượng nào không đổi trong quá trình đẳng nhiệt của một lượng khí xác định? — phát biểu Áp suất là đúng.}
{\True Hệ thức nào là định luật Boyle? — phát biểu đúng là: $p_1V_1=p_2V_2$.}
{Hệ thức nào là định luật Boyle? — phát biểu $V_1/T_1=V_2/T_2$ là đúng.}
\loigiai{
\begin{itemchoice}
\item (Đúng) Đại lượng nào không đổi trong quá trình đẳng nhiệt của một lượng khí xác định? — phát biểu đúng là: Nhiệt độ. (Vì): Trong quá trình đẳng nhiệt, nhiệt độ của lượng khí xác định là không đổi
\item (Sai) Đại lượng nào không đổi trong quá trình đẳng nhiệt của một lượng khí xác định? — phát biểu Áp suất là đúng. (Vì): Trong quá trình đẳng nhiệt, áp suất của lượng khí xác định không phải lúc nào cũng không đổi; nó thay đổi khi thể tích thay đổi theo định luật Boyle
\item (Đúng) Hệ thức nào là định luật Boyle? — phát biểu đúng là: $p_1V_1=p_2V_2$. (Vì): ịnh luật Boyle phát biểu rằng tích của áp suất và thể tích của một lượng khí xác định là không đổi ở nhiệt độ không đổi, tức là $p_1V_1=p_2V_2$
\item (Sai) Hệ thức nào là định luật Boyle? — phát biểu $V_1/T_1=V_2/T_2$ là đúng. (Vì): Hệ thức $V_1/T_1=V_2/T_2$ mô tả định luật Charles (hoặc Gay-Lussac thứ nhất), áp dụng cho quá trình đẳng áp, không phải định luật Boyle
\end{itemchoice}
}
\end{ex}

% ===== Câu 122 =====
% ID: L12C2B9-31-TL
% Mức: VD
\dangbt{Ứng dụng định luật Boyle trong bơm, xi lanh và đời sống}
\begin{bt}
% ID: L12C2B9-31-TL
% Mức: VD
Trong dạng “Ứng dụng định luật Boyle trong bơm, xi lanh và đời sống”, Thể tích giảm từ 18 lít còn 6 lít ở nhiệt độ không đổi. Áp suất tăng bao nhiêu lần?
\loigiai{
$p_2/p_1=V_1/V_2=3$.
}
\end{bt}

% ===== Câu 123 =====
% ID: L12C2B9-31-TLN
% Mức: TH
\begin{ex}
% ID: L12C2B9-31-TLN
% Mức: TH
Trong dạng “Ứng dụng định luật Boyle trong bơm, xi lanh và đời sống”, Khí 32 lít ở $150\,\text{kPa}$ được nén đẳng nhiệt còn 8 lít. Tính áp suất sau theo kPa.
\shortans{600}
\loigiai{
$p_2=150\cdot32/8=600$ kPa.
}
\end{ex}

% ===== Câu 124 =====
% ID: L12C2B9-31-TN
% Mức: VDC
\dangbt{Nhận biết quá trình đẳng nhiệt và định luật Boyle}
\begin{ex}
% ID: L12C2B9-31-TN
% Mức: VDC
Trong dạng “Nhận biết quá trình đẳng nhiệt và định luật Boyle”, Áp suất tuyệt đối của cột khí bị giam bằng
\choice
{\True áp suất khí quyển cộng hoặc trừ áp suất cột chất lỏng tùy bố trí}
{luôn bằng áp suất khí quyển}
{chỉ bằng áp suất cột chất lỏng}
{bằng không}
\loigiai{
Đáp án A. Cân bằng áp suất phải xét cả khí quyển và chênh áp thủy tĩnh.
}
\end{ex}

% ===== Câu 125 =====
% ID: L12C2B9-32-DS
% Mức: VD
\dangbt{Ứng dụng định luật Boyle trong bơm, xi lanh và đời sống}
\begin{ex}
% ID: L12C2B9-32-DS
% Mức: VD
Xét các phát biểu sau trong dạng “Ứng dụng định luật Boyle trong bơm, xi lanh và đời sống”.
\choiceTF
{\True Khi dùng định luật Boyle phải sử dụng — phát biểu đúng là: áp suất tuyệt đối.}
{Khi dùng định luật Boyle phải sử dụng — phát biểu áp suất dư là đúng.}
{\True Muốn nén khí gần đẳng nhiệt nên thực hiện — phát biểu đúng là: chậm để khí trao đổi nhiệt.}
{Muốn nén khí gần đẳng nhiệt nên thực hiện — phát biểu ở thể tích không đổi là đúng.}
\loigiai{
\begin{itemchoice}
\item (Đúng) Khi dùng định luật Boyle phải sử dụng — phát biểu đúng là: áp suất tuyệt đối. (Vì): Phương trình trạng thái dùng áp suất tuyệt đối của khí. b) Sai: trái với hệ thức hoặc điều kiện đã nêu. c) Đúng: Nén chậm giúp nhiệt độ khí gần như không đổi. d) Sai: trái với hệ thức hoặc điều kiện đã nêu
\item (Sai) Khi dùng định luật Boyle phải sử dụng — phát biểu áp suất dư là đúng.
\item (Đúng) Muốn nén khí gần đẳng nhiệt nên thực hiện — phát biểu đúng là: chậm để khí trao đổi nhiệt.
\item (Sai) Muốn nén khí gần đẳng nhiệt nên thực hiện — phát biểu ở thể tích không đổi là đúng.
\end{itemchoice}
}
\end{ex}

% ===== Câu 126 =====
% ID: L12C2B9-32-TL
% Mức: VD
\begin{bt}
% ID: L12C2B9-32-TL
% Mức: VD
Trong dạng “Ứng dụng định luật Boyle trong bơm, xi lanh và đời sống”, 120
\loigiai{
$p=960/8=120$ kPa.
}
\end{bt}

% ===== Câu 127 =====
% ID: L12C2B9-32-TLN
% Mức: TH
\begin{ex}
% ID: L12C2B9-32-TLN
% Mức: TH
Trong dạng “Ứng dụng định luật Boyle trong bơm, xi lanh và đời sống”, Khí 4 lít ở $200\,\text{kPa}$ giãn đẳng nhiệt đến áp suất $100\,\text{kPa}$. Tính thể tích sau theo lít.
\shortans{8}
\loigiai{
$V_2=200\cdot4/100=8$ lít.
}
\end{ex}

% ===== Câu 128 =====
% ID: L12C2B9-32-TN
% Mức: NB
\dangbt{Tính áp suất, thể tích trong quá trình đẳng nhiệt}
\begin{ex}
% ID: L12C2B9-32-TN
% Mức: NB
Trong dạng “Tính áp suất, thể tích trong quá trình đẳng nhiệt”, Trong quá trình đẳng nhiệt của một lượng khí xác định
\choice
{\True $pV$ không đổi}
{$p/T$ không đổi}
{$V/T$ không đổi}
{$pV/T$ tăng}
\loigiai{
Đáp án A. Định luật Boyle cho $pV=\text{hằng số}$ khi $T$ không đổi.
}
\end{ex}

% ===== Câu 129 =====
% ID: L12C2B9-33-DS
% Mức: VDC
\dangbt{Ứng dụng định luật Boyle trong bơm, xi lanh và đời sống}
\begin{ex}
% ID: L12C2B9-33-DS
% Mức: VDC
Xét các phát biểu sau trong dạng “Ứng dụng định luật Boyle trong bơm, xi lanh và đời sống”.
\choiceTF
{\True Với lượng khí cố định đẳng nhiệt, khối lượng riêng tỉ lệ — phát biểu đúng là: thuận với áp suất.}
{Với lượng khí cố định đẳng nhiệt, khối lượng riêng tỉ lệ — phát biểu nghịch với áp suất là đúng.}
{\True Trong ống có cột chất lỏng, áp suất khí bị giam được xác định bằng — phát biểu đúng là: cân bằng áp suất tại cùng độ cao.}
{Trong ống có cột chất lỏng, áp suất khí bị giam được xác định bằng — phát biểu chỉ $\rho gh$ là đúng.}
\loigiai{
\begin{itemchoice}
\item (Đúng) Với lượng khí cố định đẳng nhiệt, khối lượng riêng tỉ lệ — phát biểu đúng là: thuận với áp suất. (Vì): Vì $\rho=m/V$ và $pV$ không đổi nên $\rho\sim p$. b) Sai: trái với hệ thức hoặc điều kiện đã nêu. c) Đúng: Dùng $p_{khí}=p_0\pm\rho gh$ tùy cách bố trí. d) Sai: trái với hệ thức hoặc điều kiện đã nêu
\item (Sai) Với lượng khí cố định đẳng nhiệt, khối lượng riêng tỉ lệ — phát biểu nghịch với áp suất là đúng.
\item (Đúng) Trong ống có cột chất lỏng, áp suất khí bị giam được xác định bằng — phát biểu đúng là: cân bằng áp suất tại cùng độ cao.
\item (Sai) Trong ống có cột chất lỏng, áp suất khí bị giam được xác định bằng — phát biểu chỉ $\rho gh$ là đúng.
\end{itemchoice}
}
\end{ex}

% ===== Câu 130 =====
% ID: L12C2B9-33-TL
% Mức: VD
\begin{bt}
% ID: L12C2B9-33-TL
% Mức: VD
Trong dạng “Ứng dụng định luật Boyle trong bơm, xi lanh và đời sống”, $\rho_2/\rho_1=p_2/p_1=2$.
\loigiai{
$\rho_2/\rho_1=p_2/p_1=2$.
}
\end{bt}

% ===== Câu 131 =====
% ID: L12C2B9-33-TLN
% Mức: VD
\begin{ex}
% ID: L12C2B9-33-TLN
% Mức: VD
Trong dạng “Ứng dụng định luật Boyle trong bơm, xi lanh và đời sống”, Thể tích giảm từ 18 lít còn 6 lít ở nhiệt độ không đổi. Áp suất tăng bao nhiêu lần?
\shortans{3}
\loigiai{
$p_2/p_1=V_1/V_2=3$.
}
\end{ex}

% ===== Câu 132 =====
% ID: L12C2B9-33-TN
% Mức: NB
\dangbt{Tính áp suất, thể tích trong quá trình đẳng nhiệt}
\begin{ex}
% ID: L12C2B9-33-TN
% Mức: NB
Trong dạng “Tính áp suất, thể tích trong quá trình đẳng nhiệt”, Hệ thức đúng của định luật Boyle là
\choice
{\True $p_1V_1=p_2V_2$}
{$p_1/T_1=p_2/T_2$}
{$V_1/T_1=V_2/T_2$}
{$p_1V_2=p_2V_1$}
\loigiai{
Đáp án A. Hai trạng thái đẳng nhiệt thỏa $p_1V_1=p_2V_2$.
}
\end{ex}

% ===== Câu 133 =====
% ID: L12C2B9-34-TL
% Mức: VDC
\dangbt{Ứng dụng định luật Boyle trong bơm, xi lanh và đời sống}
\begin{bt}
% ID: L12C2B9-34-TL
% Mức: VDC
Trong dạng “Ứng dụng định luật Boyle trong bơm, xi lanh và đời sống”, undefined
\loigiai{
$V_2=140\cdot18/(420)=6$ lít.
}
\end{bt}

% ===== Câu 134 =====
% ID: L12C2B9-34-TLN
% Mức: VD
\begin{ex}
% ID: L12C2B9-34-TLN
% Mức: VD
Trong dạng “Ứng dụng định luật Boyle trong bơm, xi lanh và đời sống”, Một lượng khí có $pV=960$ kPa.lít. Khi V=8 lít, tính p theo kPa.
\shortans{120}
\loigiai{
$p=960/8=120$ kPa.
}
\end{ex}

% ===== Câu 135 =====
% ID: L12C2B9-34-TN
% Mức: TH
\dangbt{Tính áp suất, thể tích trong quá trình đẳng nhiệt}
\begin{ex}
% ID: L12C2B9-34-TN
% Mức: TH
Trong dạng “Tính áp suất, thể tích trong quá trình đẳng nhiệt”, Nén đẳng nhiệt để thể tích giảm một nửa thì áp suất
\choice
{\True tăng gấp đôi}
{giảm một nửa}
{không đổi}
{tăng bốn lần}
\loigiai{
Đáp án A. p tỉ lệ nghịch $V$ nên $V$ giảm 2 lần thì p tăng 2 lần.
}
\end{ex}

% ===== Câu 136 =====
% ID: L12C2B9-35-TLN
% Mức: VD
\dangbt{Ứng dụng định luật Boyle trong bơm, xi lanh và đời sống}
\begin{ex}
% ID: L12C2B9-35-TLN
% Mức: VD
Trong dạng “Ứng dụng định luật Boyle trong bơm, xi lanh và đời sống”, Áp suất đẳng nhiệt tăng từ $130\,\text{kPa}$ lên $260\,\text{kPa}$. Khối lượng riêng tăng bao nhiêu lần?
\shortans{2}
\loigiai{
$\rho_2/\rho_1=p_2/p_1=2$.
}
\end{ex}

% ===== Câu 137 =====
% ID: L12C2B9-35-TN
% Mức: TH
\dangbt{Tính áp suất, thể tích trong quá trình đẳng nhiệt}
\begin{ex}
% ID: L12C2B9-35-TN
% Mức: TH
Trong dạng “Tính áp suất, thể tích trong quá trình đẳng nhiệt”, Đường đẳng nhiệt trong hệ p- $V$ là
\choice
{\True nhánh hypebol}
{đường thẳng qua gốc}
{đường tròn}
{parabol}
\loigiai{
Đáp án A. Phương trình $p=C/V$ là một nhánh hypebol.
}
\end{ex}

% ===== Câu 138 =====
% ID: L12C2B9-36-TLN
% Mức: VDC
\dangbt{Ứng dụng định luật Boyle trong bơm, xi lanh và đời sống}
\begin{ex}
% ID: L12C2B9-36-TLN
% Mức: VDC
Trong dạng “Ứng dụng định luật Boyle trong bơm, xi lanh và đời sống”, Khí 18 lít ở $140\,\text{kPa}$ giãn/nén đẳng nhiệt đến $420\,\text{kPa}$. Tính thể tích theo lít.
\shortans{6}
\loigiai{
$V_2=140\cdot18/(420)=6$ lít.
}
\end{ex}

% ===== Câu 139 =====
% ID: L12C2B9-36-TN
% Mức: TH
\dangbt{Tính áp suất, thể tích trong quá trình đẳng nhiệt}
\begin{ex}
% ID: L12C2B9-36-TN
% Mức: TH
Trong dạng “Tính áp suất, thể tích trong quá trình đẳng nhiệt”, Trên đồ thị p- $V$, đường đẳng nhiệt có nhiệt độ cao hơn nằm
\choice
{\True xa gốc tọa độ hơn}
{gần gốc hơn}
{trùng trục $V$}
{trùng trục p}
\loigiai{
Đáp án A. Với cùng $V$, $T$ cao hơn cho p lớn hơn nên đường nằm xa gốc.
}
\end{ex}

% ===== Câu 140 =====
% ID: L12C2B9-37-TN
% Mức: VD
\begin{ex}
% ID: L12C2B9-37-TN
% Mức: VD
Trong dạng “Tính áp suất, thể tích trong quá trình đẳng nhiệt”, Khí có $p_1=1$ atm, $V_1=4$ lít, nén còn 2 lít. $p_2$ bằng
\choice
{\True 2 atm}
{0,5 atm}
{4 atm}
{8 atm}
\loigiai{
Đáp án A. $p_2=p_1V_1/V_2=2$ atm.
}
\end{ex}

% ===== Câu 141 =====
% ID: L12C2B9-38-TN
% Mức: VD
\begin{ex}
% ID: L12C2B9-38-TN
% Mức: VD
Trong dạng “Tính áp suất, thể tích trong quá trình đẳng nhiệt”, Trong quá trình đẳng nhiệt, khối lượng riêng của khí tỉ lệ
\choice
{\True thuận với áp suất}
{nghịch với áp suất}
{thuận với thể tích}
{không phụ thuộc áp suất}
\loigiai{
Đáp án A. Với khối lượng cố định, $\rho=m/V$ và pV không đổi nên $\rho\sim p$.
}
\end{ex}

% ===== Câu 142 =====
% ID: L12C2B9-39-TN
% Mức: VD
\begin{ex}
% ID: L12C2B9-39-TN
% Mức: VD
Trong dạng “Tính áp suất, thể tích trong quá trình đẳng nhiệt”, Bơm xe đạp nóng lên rõ rệt thì quá trình nén
\choice
{\True không hoàn toàn đẳng nhiệt}
{luôn đẳng nhiệt chính xác}
{đẳng áp}
{đẳng tích}
\loigiai{
Đáp án A. Nhiệt độ thay đổi nên không thể áp dụng Boyle chính xác như quá trình đẳng nhiệt.
}
\end{ex}

% ===== Câu 143 =====
% ID: L12C2B9-40-TN
% Mức: VDC
\begin{ex}
% ID: L12C2B9-40-TN
% Mức: VDC
Trong dạng “Tính áp suất, thể tích trong quá trình đẳng nhiệt”, Để quá trình gần đẳng nhiệt nên nén khí
\choice
{\True chậm để trao đổi nhiệt}
{thật nhanh}
{trong bình cách nhiệt}
{với thể tích không đổi}
\loigiai{
Đáp án A. Nén chậm giúp khí trao đổi nhiệt và giữ nhiệt độ gần không đổi.
}
\end{ex}

% ===== Câu 144 =====
% ID: L12C2B9-41-TN
% Mức: VDC
\begin{ex}
% ID: L12C2B9-41-TN
% Mức: VDC
Trong dạng “Tính áp suất, thể tích trong quá trình đẳng nhiệt”, Áp suất tuyệt đối của cột khí bị giam bằng
\choice
{\True áp suất khí quyển cộng hoặc trừ áp suất cột chất lỏng tùy bố trí}
{luôn bằng áp suất khí quyển}
{chỉ bằng áp suất cột chất lỏng}
{bằng không}
\loigiai{
Đáp án A. Cân bằng áp suất phải xét cả khí quyển và chênh áp thủy tĩnh.
}
\end{ex}

% ===== Câu 145 =====
% ID: L12C2B9-42-TN
% Mức: NB
\dangbt{Đồ thị đẳng nhiệt trên các hệ tọa độ}
\begin{ex}
% ID: L12C2B9-42-TN
% Mức: NB
Trong dạng “Đồ thị đẳng nhiệt trên các hệ tọa độ”, Trong quá trình đẳng nhiệt của một lượng khí xác định
\choice
{\True $pV$ không đổi}
{$p/T$ không đổi}
{$V/T$ không đổi}
{$pV/T$ tăng}
\loigiai{
Đáp án A. Định luật Boyle cho $pV=\text{hằng số}$ khi $T$ không đổi.
}
\end{ex}

% ===== Câu 146 =====
% ID: L12C2B9-43-TN
% Mức: NB
\begin{ex}
% ID: L12C2B9-43-TN
% Mức: NB
Trong dạng “Đồ thị đẳng nhiệt trên các hệ tọa độ”, Hệ thức đúng của định luật Boyle là
\choice
{\True $p_1V_1=p_2V_2$}
{$p_1/T_1=p_2/T_2$}
{$V_1/T_1=V_2/T_2$}
{$p_1V_2=p_2V_1$}
\loigiai{
Đáp án A. Hai trạng thái đẳng nhiệt thỏa $p_1V_1=p_2V_2$.
}
\end{ex}

% ===== Câu 147 =====
% ID: L12C2B9-44-TN
% Mức: TH
\begin{ex}
% ID: L12C2B9-44-TN
% Mức: TH
Trong dạng “Đồ thị đẳng nhiệt trên các hệ tọa độ”, Nén đẳng nhiệt để thể tích giảm một nửa thì áp suất
\choice
{\True tăng gấp đôi}
{giảm một nửa}
{không đổi}
{tăng bốn lần}
\loigiai{
Đáp án A. p tỉ lệ nghịch $V$ nên $V$ giảm 2 lần thì p tăng 2 lần.
}
\end{ex}

% ===== Câu 148 =====
% ID: L12C2B9-45-TN
% Mức: TH
\begin{ex}
% ID: L12C2B9-45-TN
% Mức: TH
Trong dạng “Đồ thị đẳng nhiệt trên các hệ tọa độ”, Đường đẳng nhiệt trong hệ p- $V$ là
\choice
{\True nhánh hypebol}
{đường thẳng qua gốc}
{đường tròn}
{parabol}
\loigiai{
Đáp án A. Phương trình $p=C/V$ là một nhánh hypebol.
}
\end{ex}

% ===== Câu 149 =====
% ID: L12C2B9-46-TN
% Mức: TH
\begin{ex}
% ID: L12C2B9-46-TN
% Mức: TH
Trong dạng “Đồ thị đẳng nhiệt trên các hệ tọa độ”, Trên đồ thị p- $V$, đường đẳng nhiệt có nhiệt độ cao hơn nằm
\choice
{\True xa gốc tọa độ hơn}
{gần gốc hơn}
{trùng trục $V$}
{trùng trục p}
\loigiai{
Đáp án A. Với cùng $V$, $T$ cao hơn cho p lớn hơn nên đường nằm xa gốc.
}
\end{ex}

% ===== Câu 150 =====
% ID: L12C2B9-47-TN
% Mức: VD
\begin{ex}
% ID: L12C2B9-47-TN
% Mức: VD
Trong dạng “Đồ thị đẳng nhiệt trên các hệ tọa độ”, Khí có $p_1=1$ atm, $V_1=4$ lít, nén còn 2 lít. $p_2$ bằng
\choice
{\True 2 atm}
{0,5 atm}
{4 atm}
{8 atm}
\loigiai{
Đáp án A. $p_2=p_1V_1/V_2=2$ atm.
}
\end{ex}

% ===== Câu 151 =====
% ID: L12C2B9-48-TN
% Mức: VD
\begin{ex}
% ID: L12C2B9-48-TN
% Mức: VD
Trong dạng “Đồ thị đẳng nhiệt trên các hệ tọa độ”, Trong quá trình đẳng nhiệt, khối lượng riêng của khí tỉ lệ
\choice
{\True thuận với áp suất}
{nghịch với áp suất}
{thuận với thể tích}
{không phụ thuộc áp suất}
\loigiai{
Đáp án A. Với khối lượng cố định, $\rho=m/V$ và pV không đổi nên $\rho\sim p$.
}
\end{ex}

% ===== Câu 152 =====
% ID: L12C2B9-49-TN
% Mức: VD
\begin{ex}
% ID: L12C2B9-49-TN
% Mức: VD
Trong dạng “Đồ thị đẳng nhiệt trên các hệ tọa độ”, Bơm xe đạp nóng lên rõ rệt thì quá trình nén
\choice
{\True không hoàn toàn đẳng nhiệt}
{luôn đẳng nhiệt chính xác}
{đẳng áp}
{đẳng tích}
\loigiai{
Đáp án A. Nhiệt độ thay đổi nên không thể áp dụng Boyle chính xác như quá trình đẳng nhiệt.
}
\end{ex}

% ===== Câu 153 =====
% ID: L12C2B9-50-TN
% Mức: VDC
\begin{ex}
% ID: L12C2B9-50-TN
% Mức: VDC
Trong dạng “Đồ thị đẳng nhiệt trên các hệ tọa độ”, Để quá trình gần đẳng nhiệt nên nén khí
\choice
{\True chậm để trao đổi nhiệt}
{thật nhanh}
{trong bình cách nhiệt}
{với thể tích không đổi}
\loigiai{
Đáp án A. Nén chậm giúp khí trao đổi nhiệt và giữ nhiệt độ gần không đổi.
}
\end{ex}

% ===== Câu 154 =====
% ID: L12C2B9-51-TN
% Mức: VDC
\begin{ex}
% ID: L12C2B9-51-TN
% Mức: VDC
Trong dạng “Đồ thị đẳng nhiệt trên các hệ tọa độ”, Áp suất tuyệt đối của cột khí bị giam bằng
\choice
{\True áp suất khí quyển cộng hoặc trừ áp suất cột chất lỏng tùy bố trí}
{luôn bằng áp suất khí quyển}
{chỉ bằng áp suất cột chất lỏng}
{bằng không}
\loigiai{
Đáp án A. Cân bằng áp suất phải xét cả khí quyển và chênh áp thủy tĩnh.
}
\end{ex}

% ===== Câu 155 =====
% ID: L12C2B9-52-TN
% Mức: NB
\dangbt{Ứng dụng định luật Boyle trong bơm, xi lanh và đời sống}
\begin{ex}
% ID: L12C2B9-52-TN
% Mức: NB
Trong dạng “Ứng dụng định luật Boyle trong bơm, xi lanh và đời sống”, Đại lượng nào không đổi trong quá trình đẳng nhiệt của một lượng khí xác định?
\choice
{\True Nhiệt độ}
{Áp suất}
{Thể tích}
{Khối lượng riêng}
\loigiai{
Đáp án A. Quá trình đẳng nhiệt có nhiệt độ không đổi.
}
\end{ex}

% ===== Câu 156 =====
% ID: L12C2B9-53-TN
% Mức: NB
\begin{ex}
% ID: L12C2B9-53-TN
% Mức: NB
Trong dạng “Ứng dụng định luật Boyle trong bơm, xi lanh và đời sống”, Hệ thức nào là định luật Boyle?
\choice
{\True $p_1V_1=p_2V_2$}
{$p_1/T_1=p_2/T_2$}
{$V_1/T_1=V_2/T_2$}
{$p_1V_2=p_2V_1$}
\loigiai{
Đáp án A. Ở nhiệt độ không đổi, $pV$ là hằng số.
}
\end{ex}

% ===== Câu 157 =====
% ID: L12C2B9-54-TN
% Mức: TH
\begin{ex}
% ID: L12C2B9-54-TN
% Mức: TH
Trong dạng “Ứng dụng định luật Boyle trong bơm, xi lanh và đời sống”, Khí có áp suất $125\,\text{kPa}$, thể tích 20 lít được nén đẳng nhiệt còn 4 lít. Áp suất sau là
\choice
{\True $625\,\text{kPa}$}
{$25\,\text{kPa}$}
{$125\,\text{kPa}$}
{$3125\,\text{kPa}$}
\loigiai{
Đáp án A. $p_2=p_1V_1/V_2=125\cdot20/4=625$ kPa.
}
\end{ex}

% ===== Câu 158 =====
% ID: L12C2B9-55-TN
% Mức: TH
\begin{ex}
% ID: L12C2B9-55-TN
% Mức: TH
Trong dạng “Ứng dụng định luật Boyle trong bơm, xi lanh và đời sống”, Khí có áp suất $130\,\text{kPa}$, thể tích 6 lít giãn đẳng nhiệt đến 12 lít. Áp suất sau là
\choice
{\True $65\,\text{kPa}$}
{$260\,\text{kPa}$}
{$130\,\text{kPa}$}
{$32,5\,\text{kPa}$}
\loigiai{
Đáp án A. $p_2=130\cdot6/12=65$ kPa.
}
\end{ex}

% ===== Câu 159 =====
% ID: L12C2B9-56-TN
% Mức: VD
\dangbt{Từ trường của dòng điện thẳng, dòng điện tròn và ống dây}
\begin{ex}
% ID: L12C2B9-56-TN
% Mức: VD
Cơ sở Vật lí đúng dùng để giải từ trường của dòng điện thẳng, dòng điện tròn và ống dây là gì?
\choice
{Có thể bỏ qua phương, chiều và đơn vị của các đại lượng trong mọi trường hợp.}
{Đại lượng đang xét không phụ thuộc các điều kiện vật lí nêu trong bài.}
{Từ trường trong lòng ống dây dài luôn bằng không.}
{\True Đường sức quanh dây dẫn thẳng là các đường tròn đồng tâm; từ trường trong lòng ống dây dài gần đều.}
\loigiai{
Đường sức quanh dây dẫn thẳng là các đường tròn đồng tâm; từ trường trong lòng ống dây dài gần đều. Vì vậy phương án $D$ đúng; các phương án còn lại trái với định nghĩa, định luật hoặc điều kiện áp dụng.
}
\end{ex}

% ===== Câu 160 =====
% ID: L12C3B14-101-TN
% Mức: VDC
\dangbt{Bài toán tổng hợp về quá trình đẳng nhiệt}
\begin{ex}
% ID: L12C3B14-101-TN
% Mức: VDC
Trong ống có cột chất lỏng, áp suất khí bị giam được xác định bằng
\choice
{\True cân bằng áp suất tại cùng độ cao}
{chỉ áp suất khí quyển}
{chỉ $\rho gh$}
{luôn bằng không}
\loigiai{
Đáp án A. Dùng $p_{khí}=p_0\pm\rho gh$ tùy cách bố trí.
}
\end{ex}
